% This LaTeX document needs to be compiled with XeLaTeX.
\documentclass[10pt]{article}
\usepackage[utf8]{inputenc}
\usepackage{amsmath}
\usepackage{amsfonts}
\usepackage{amssymb}
\usepackage[version=4]{mhchem}
\usepackage{stmaryrd}
\usepackage{graphicx}
\usepackage[export]{adjustbox}
\graphicspath{ {./images/} }
\usepackage{bbold}
\usepackage{caption}
\usepackage{hyperref}
\hypersetup{colorlinks=true, linkcolor=blue, filecolor=magenta, urlcolor=cyan,}
\urlstyle{same}
\usepackage[fallback]{xeCJK}
\usepackage{polyglossia}
\usepackage{fontspec}
\usepackage{newunicodechar}
\IfFontExistsTF{Noto Serif CJK JP}
{\setCJKmainfont{Noto Serif CJK JP}}
{\IfFontExistsTF{STSong}
  {\setCJKmainfont{STSong}}
  {\IfFontExistsTF{Droid Sans Fallback}
    {\setCJKmainfont{Droid Sans Fallback}}
    {\setCJKmainfont{SimSun}}
}}

\setmainlanguage{english}
\IfFontExistsTF{CMU Serif}
{\setmainfont{CMU Serif}}
{\IfFontExistsTF{DejaVu Sans}
  {\setmainfont{DejaVu Sans}}
  {\setmainfont{Georgia}}
}

%New command to display footnote whose markers will always be hidden
\let\svthefootnote\thefootnote
\newcommand\blfootnotetext[1]{%
  \let\thefootnote\relax\footnote{#1}%
  \addtocounter{footnote}{-1}%
  \let\thefootnote\svthefootnote%
}

%Overriding the \footnotetext command to hide the marker if its value is `0`
\let\svfootnotetext\footnotetext
\renewcommand\footnotetext[2][?]{%
  \if\relax#1\relax%
    \ifnum\value{footnote}=0\blfootnotetext{#2}\else\svfootnotetext{#2}\fi%
  \else%
    \if?#1\ifnum\value{footnote}=0\blfootnotetext{#2}\else\svfootnotetext{#2}\fi%
    \else\svfootnotetext[#1]{#2}\fi%
  \fi
}

\newunicodechar{↪}{\ifmmode\hookrightarrow\else{$\hookrightarrow$}\fi}
\newunicodechar{→}{\ifmmode\rightarrow\else{$\rightarrow$}\fi}
\newunicodechar{⇒}{\ifmmode\Rightarrow\else{$\Rightarrow$}\fi}
\newunicodechar{⇔}{\ifmmode\Leftrightarrow\else{$\Leftrightarrow$}\fi}

\begin{document}
\captionsetup{singlelinecheck=false}
\section*{CHAPTER 1}
\section*{$\infty$-Cosmoi and Their Homotopy 2-Categories}
In this chapter, we introduce a framework to develop the formal category theory of $\infty$-categories, which goes by the name of an $\infty$-cosmos. Informally, an $\infty$-cosmos is an ( $\infty, 2$ )-category - a category enriched over $(\infty, 1)$-categories - that is equipped with $(\infty, 2)$-categorical limits. In the motivating examples of $\infty$-cosmoi, the objects are $\infty$-categories in some model. To focus this abstract theory on its intended interpretation, we recast " $\infty$-category" as a technical term, reserved to mean an object of some $\infty$-cosmos.

Unexpectedly, the motivating examples permit us to use a quite strict interpretation of "( $\infty, 2$ )-category with $(\infty, 2)$-categorical limits": an $\infty$-cosmos is a particular type of simplicially enriched category and the $(\infty, 2)$-categorical limits are modeled by simplicially enriched limits. More precisely, an $\infty$-cosmos is a category enriched over quasi-categories, these being one of the models of $(\infty, 1)$-cate-gories defined as certain simplicial sets. The $(\infty, 2)$-categorical limits are defined as limits of diagrams involving specified maps called isofibrations, which have no intrinsic homotopical meaning - since any functor between $\infty$-categories is equivalent to an isofibration - but allow us to consider strictly commuting diagrams.

In §1.1, we introduce quasi-categories, reviewing the classical results that are needed to show that quasi-categories themselves assemble into an $\infty$-cosmos - the prototypical example. General $\infty$ -cosmoi are defined in §1.2, where several examples are given and their basic properties are established. In §1.3, we turn our attention to cosmological functors between $\infty$-cosmoi, which preserve all of the defining structure. Cosmological functors serve dual purposes, on the one hand providing technical simplifications in many proofs, and then later on serving as the "change of model" functors that establish the model independence of $\infty$-category theory.

Finally, in §1.4, we introduce a strict 2-category whose objects are $\infty$-categories, whose 1-cells are the $\infty$-functors between them, and whose 2-cells define $\infty$-natural transformations between these. Any $\infty$-cosmos has a 2-category of this sort, which we refer to as the homotopy 2-category of the $\infty$-cosmos. In fact, the reader who is eager to get on to the development of the category theory of $\infty$-categories can skip this chapter on first reading, taking the existence of the homotopy 2-category for granted, and start with Chapter 2.

\subsection*{1.1. Quasi-Categories}
Before introducing an axiomatic framework that allows us to develop $\infty$-category theory in general, we first consider one model in particular: quasi-categories, which were introduced in 1973 by Boardman and Vogt [21] in their study of homotopy coherent diagrams. Ordinary 1-categories give examples of quasi-categories via the construction of Definition 1.1.4. Joyal first undertook the task of extending 1-category theory to quasi-category theory in [61] and [63] and in several unpublished draft book manuscripts. The majority of the results in this section are due to him.\\
1.1.1. Notation (the simplex category). Let $\boldsymbol{\Delta}$ denote the simplex category of finite nonempty ordinals $[n]=\{0<1<\cdots<n\}$ and order-preserving maps. These include in particular the

$$
\begin{array}{ccc}
\text { elementary face operators } & {[n-1] \stackrel{\delta^{i}}{\rightarrow}[n]} & 0 \leq i \leq n \\
\text { elementary degeneracy operators } & {[n+1] \xrightarrow{\sigma^{i}}[n]} & 0 \leq i \leq n
\end{array}
$$

whose images, respectively, omit and double up on the element $i \in[n]$. Every morphism in $\boldsymbol{\Delta}$ factors uniquely as an epimorphism followed by a monomorphism; these epimorphisms, the degeneracy operators, decompose as composites of elementary degeneracy operators, while the monomorphisms, the face operators, decompose as composites of elementary face operators.

The category of simplicial sets is the category s et $:=$ Set ${ }^{\Delta^{\text {op }}}$ of presheaves on the simplex category. We write $\Delta[n]$ for the standard $n$-simplex the simplicial set represented by $[n] \in \Delta$, and $\Lambda^{k}[n] \subset$ $\partial \Delta[n] \subset \Delta[n]$ for its $k$-horn and boundary sphere, respectively. The sphere $\partial \Delta[n]$ is the simplicial subset generated by the codimension-one faces of the $n$-simplex, while the horn $\Lambda^{k}[n]$ is the further simplicial subset that omits the face opposite the vertex $k$.

Given a simplicial set $X$, it is conventional to write $X_{n}$ for the set of $n$-simplices, defined by evaluating at $[n] \in \boldsymbol{\Delta}$. By the Yoneda lemma, each $n$-simplex $x \in X_{n}$ corresponds to a map of simplicial sets $x: \Delta[n] \rightarrow X$. Accordingly, we write $x \cdot \delta^{i}$ for the $i$ th face of the $n$-simplex, an $(n-1)$-simplex classified by the composite map

$$
\Delta[n-1] \xrightarrow{\delta^{i}} \Delta[n] \xrightarrow{x} X .
$$

The right action of the face operator defines a map $X_{n} \xrightarrow{. \delta^{i}} X_{n-1}$. Geometrically, $x \cdot \delta^{i}$ is the "face opposite the vertex $i$ " in the $n$-simplex $x$.\\
1.1.2. Definition (quasi-category). A quasi-category is a simplicial set $A$ in which any inner horn can be extended to a simplex, solving the displayed lifting problem:\\
\includegraphics[max width=\textwidth, alt={}, center]{b375083c-8174-49a1-a017-e977744038d8-02_198_731_1579_703}

(1.1.3)

Quasi-categories were first introduced by Boardman and Vogt [21] under the name "weak Kan complexes," a Kan complex being a simplicial set admitting extensions as in (1.1.3) along all horn inclusions $n \geq 1,0 \leq k \leq n$. Since any topological space can be encoded as a Kan complex, ${ }^{1}$ in this way spaces provide examples of quasi-categories.

Categories also provide examples of quasi-categories via the nerve construction.\\
1.1.4. Definition (nerve). The category $\boldsymbol{C} a t$ of 1 -categories embeds fully faithfully into the category of simplicial sets via the nerve functor. An $n$-simplex in the nerve of a 1 -category $C$ is a sequence of $n$ composable arrows in $C$, or equally a functor $\mathrm{m}+\mathbb{1} \rightarrow C$ from the ordinal category $\mathrm{m}+\mathbb{1}:=[n]$ with objects $0, \ldots, n$ and a unique arrow $i \rightarrow j$ just when $i \leq j$.

\footnotetext{${ }^{1}$ The total singular complex construction defines a functor from topological spaces to simplicial sets that is an equivalence on their respective homotopy categories - weak homotopy types of spaces correspond to homotopy equivalence classes of Kan complexes [93, §II.2]. The left adjoint constructed by Exercise 1.1.i "geometrically realizes" a simplicial set as a topological space.
}The map $[n] \mapsto \mathrm{m}+\mathbb{1}$ defines a fully faithful embedding $\boldsymbol{\Delta} \hookrightarrow$ Cat. From this point of view, the nerve functor can be described as a "restricted Yoneda embedding" which carries a category $C$ to the restriction of the representable functor hom(-, $C$ ) to the image of this inclusion. More general "nerve-type constructions" are described in Exercise 1.1.i.\\
1.1.5. Remark. The nerve of a category $C$ is 2-coskeletal as a simplicial set, meaning that every sphere $\partial \Delta[n] \rightarrow C$ with $n \geq 3$ is filled uniquely by an $n$-simplex in $C$ (see Definition C.5.2). Note a sphere $\partial \Delta[2] \rightarrow C$ extends to a 2-simplex if and only if that arrow along its diagonal edge is the composite of the arrows along the edges in the inner horn $\Lambda^{1}[2] \subset \partial \Delta[2] \rightarrow C$. The simplices in dimension 3 and above witness the associativity of the composition of the path of composable arrows found along their spine, the 1-skeletal simplicial subset formed by the edges connecting adjacent vertices. In fact, as suggested by the proof of Proposition 1.1.6, any simplicial set in which inner horns admit unique fillers is isomorphic to the nerve of a 1-category (see Exercise 1.1.iv).

We decline to introduce explicit notation for the nerve functor, preferring instead to identify 1 -categories with their nerves. As we shall discover the theory of 1 -categories extends to $\infty$-categories modeled as quasi-categories in such a way that the restriction of each $\infty$-categorical concept along the nerve embedding recovers the corresponding 1-categorical concept. For instance, the standard simplex $\Delta[n]$ is isomorphic to the nerve of the ordinal category $\mathrm{m}+1$, and we frequently adopt the latter notation - writing $\mathbb{1}:=\Delta[0], 2:=\Delta[1], 3:=\Delta[2]$, and so on - to suggest the correct categorical intuition.

To begin down this path, we must first verify the implicit assertion that has just been made:

\subsection*{1.1.6. Proposition (nerves are quasi-categories). Nerves of categories are quasi-categories.}
Proof. Via the isomorphism $C \cong \operatorname{cosk}_{2} C$ from Remark 1.1.5 and the adjunction $\operatorname{sk}_{2} \dashv \operatorname{cosk}_{2}$ of C.5.2, the required lifting problem displayed below-left transposes to the one displayed below-right:\\
\includegraphics[max width=\textwidth, alt={}, center]{b375083c-8174-49a1-a017-e977744038d8-03_203_1027_1504_548}\\
The functor $\mathrm{sk}_{2}$ replaces a simplicial set by its 2-skeleton, the simplicial subset generated by the simplices of dimension at most two. For $n \geq 4$, the inclusion $\operatorname{sk}_{2} \Lambda^{k}[n] \hookrightarrow \operatorname{sk}_{2} \Delta[n]$ is an isomorphism, in which case the lifting problems on the right admit (unique) solutions. So it remains only to solve the lifting problems on the left in the cases $n=2$ and $n=3$.

To that end consider\\
\includegraphics[max width=\textwidth, alt={}, center]{b375083c-8174-49a1-a017-e977744038d8-03_196_1035_1991_544}\\
An inner horn $\Lambda^{1}[2] \rightarrow C$ defines a composable pair of arrows in $C$; an extension to a 2 -simplex exists precisely because any composable pair of arrows admits a (unique) composite.

An inner horn $\Lambda^{1}[3] \rightarrow C$ specifies the data of three composable arrows in $C$, as displayed in the following diagram, together with the composites $g f, h g$, and $(h g) f$.\\
\includegraphics[max width=\textwidth, alt={}, center]{b375083c-8174-49a1-a017-e977744038d8-04_304_372_373_874}\\
Because composition is associative, the arrow $(h g) f$ is also the composite of $g f$ followed by $h$, which proves that the 2 -simplex opposite the vertex $c_{1}$ is present in $C$; by 2-coskeletality, the 3 -simplex filling this boundary sphere is also present in $C$. The filler for a horn $\Lambda^{2}[3] \rightarrow C$ is constructed similarly. \(\square\)\\
1.1.7. Definition (homotopy relation on 1 -simplices). A parallel pair of 1 -simplices $f, g$ in a simplicial set $X$ are homotopic if there exists a 2-simplex whose boundary takes either of the following forms ${ }^{2}$\\
\includegraphics[max width=\textwidth, alt={}, center]{b375083c-8174-49a1-a017-e977744038d8-04_151_239_986_763}\\
\includegraphics[max width=\textwidth, alt={}, center]{b375083c-8174-49a1-a017-e977744038d8-04_151_237_986_1122}

(1.1.8)\\
or if $f$ and $g$ are in the same equivalence class generated by this relation.\\
In a quasi-category, the relation witnessed by either of the types of 2-simplex on display in (1.1.8) is an equivalence relation and these equivalence relations coincide.\\
1.1.9. Lemma (homotopic 1-simplices in a quasi-category). Parallel 1-simplices $f$ and $g$ in a quasi-category are homotopic if and only if there exists a 2-simplex of any or equivalently all of the forms displayed in (1.1.8).

Proof. Exercise 1.1.ii. \(\square\)\\[0pt]
1.1.10. Definition (the homotopy category [44, §2.4]). By 1-truncating, any simplicial set $X$ has an underlying reflexive directed graph with the 0 -simplices of $X$ defining the objects and the 1 -simplices defining the arrows:

$$
X_{1} \xrightarrow[\cdot \delta^{0}]{\xrightarrow{\stackrel{\cdot \delta^{1}}{\leftarrow \cdot \sigma^{0}}-}} X_{0},
$$

By convention, the source of an arrow $f \in X_{1}$ is its 0th face $f \cdot \delta^{1}$ (the face opposite 1) while the target is its 1st face $f \cdot \delta^{0}$ (the face opposite 0). The free category on this reflexive directed graph has $X_{0}$ as its object set, degenerate 1 -simplices serving as identity morphisms, and nonidentity morphisms defined to be finite directed paths of nondegenerate 1-simplices. The homotopy category $\mathrm{h} X$ of $X$ is the quotient of the free category on its underlying reflexive directed graph by the congruence ${ }^{3}$ generated by imposing a composition relation $h=g \circ f$ witnessed by 2-simplices\\
\includegraphics[max width=\textwidth, alt={}, center]{b375083c-8174-49a1-a017-e977744038d8-04_146_285_2144_918}

\footnotetext{${ }^{2}$ The symbol " = " is used in diagrams to denote a degenerate simplex or an identity arrow.\\
${ }^{3}$ A binary relation \~{} on parallel arrows of a 1-category is a congruence if it is an equivalence relation that is closed under pre- and post-composition: if $f \sim g$ then $h f k \sim h g k$.
}This relation implies in particular that homotopic 1-simplices represent the same arrow in the homotopy category.

The homotopy category of the nerve of a 1-category is isomorphic to the original category, as the 2-simplices in the nerve witness all of the composition relations satisfied by the arrows in the underlying reflexive directed graph. Indeed, the natural isomorphism $\mathrm{h} C \cong C$ forms the counit of an adjunction, embedding Cat as a reflective subcategory of sSet.\\
1.1.11. Proposition. The nerve embedding admits a left adjoint, namely the functor which sends a simplicial set to its homotopy category:\\
\includegraphics[max width=\textwidth, alt={}, center]{b375083c-8174-49a1-a017-e977744038d8-05_118_382_697_870}

The adjunction of Proposition 1.1.11 exists for formal reasons (see Exercise 1.1.i), but nevertheless, a direct proof can be enlightening.

Proof. For any simplicial set $X$, there is a natural map from $X$ to the nerve of its homotopy category $\mathrm{h} X$; since nerves are 2 -coskeletal, it suffices to define the map $\mathrm{sk}_{2} X \rightarrow \mathrm{~h} X$, and this is given immediately by the construction of Definition 1.1.10. Note that the quotient map $X \rightarrow \mathrm{~h} X$ becomes an isomorphism upon applying the homotopy category functor and is already an isomorphism whenever $X$ is the nerve of a category. Thus the adjointness follows from Lemma B.4.2 or by direct verification of the triangle equalities. \(\square\)

The homotopy category of a quasi-category admits a simplified description.\\
1.1.12. Lemma (the homotopy category of a quasi-category). If $A$ is a quasi-category then its homotopy category $\mathrm{h} A$ has

\begin{itemize}
  \item[-] the set of 0 -simplices $A_{0}$ as its objects
  \item[-] the set of homotopy classes of 1 -simplices $A_{1}$ as its arrows
  \item[-] the identity arrow at $a \in A_{0}$ represented by the degenerate 1-simplex $a \cdot \sigma^{0} \in A_{1}$
  \item[-] a composition relation $h=g \circ f$ in $\mathrm{h} A$ between the homotopy classes of arrows represented by any given 1 -simplices $f, g, h \in A_{1}$ if and only if there exists a 2 -simplex with boundary
\end{itemize}

\begin{figure}[h]
\begin{center}
  \includegraphics[alt={},max width=\textwidth]{b375083c-8174-49a1-a017-e977744038d8-05_148_284_1723_919}
\captionsetup{labelformat=empty}
\caption{Proof. Exercise 1.1.iii.}
\end{center}
\end{figure}

1.1.13. Definition (isomorphism in a quasi-category). A 1-simplex in a quasi-category is an isomorph$\mathbf{i s m}^{4}$ just when it represents an isomorphism in the homotopy category. By Lemma 1.1.12 this means that $f: a \rightarrow b$ is an isomorphism if and only if there exists a 1 -simplex $f^{-1}: b \rightarrow a$ together with a pair of 2-simplices\\
\includegraphics[max width=\textwidth, alt={}, center]{b375083c-8174-49a1-a017-e977744038d8-05_139_234_2178_768}\\
\includegraphics[max width=\textwidth, alt={}, center]{b375083c-8174-49a1-a017-e977744038d8-05_130_232_2187_1122}

\footnotetext{${ }^{4}$ Joyal refers to these maps as "isomorphisms" while Lurie refers to them as "equivalences." We prefer, wherever possible, to use the same term for $\infty$-categorical concepts as for the analogous 1-categorical ones.
}The properties of the isomorphisms in a quasi-category are most easily proved by arguing in a closely related category where simplicial sets have the additional structure of a "marking" on a specified subset of the 1-simplices; maps of these so-called marked simplicial sets must then preserve the markings (see Definition D.1.1). For instance, each quasi-category has a natural marking, where the marked 1-simplices are exactly the isomorphisms (see Definition D.4.5). Since the property of being an isomorphism in a quasi-category is witnessed by the presence of 2-simplices with a particular boundary, every map between quasi-categories preserves isomorphisms, inducing a map of the corresponding naturally marked quasi-categories. Because marked simplicial sets seldom appear outside of the proofs of certain combinatorial lemmas about the isomorphisms in quasi-categories, we save the details for Appendix D.

Let us now motivate the first of several results proven using marked techniques. A quasi-category $A$ is defined to have extensions along all inner horns. But when the initial or final edges, respectively, of an outer horn $\Lambda^{0}[2] \rightarrow A$ or $\Lambda^{2}[2] \rightarrow A$ map to isomorphisms in $A$, then a filler\\
\includegraphics[max width=\textwidth, alt={}, center]{b375083c-8174-49a1-a017-e977744038d8-06_147_283_858_717}\\
\includegraphics[max width=\textwidth, alt={}, center]{b375083c-8174-49a1-a017-e977744038d8-06_146_282_859_1123}\\
should intuitively exist. The higher-dimensional "special outer horns" behave similarly:\\
1.1.14. Proposition (special outer horn filling). Any quasi-category $A$ admits fillers for those outer horns\\
\includegraphics[max width=\textwidth, alt={}, center]{b375083c-8174-49a1-a017-e977744038d8-06_205_1013_1167_563}\\
in which the edges $\left.g\right|_{\{0,1\}}$ and $\left.h\right|_{\{n-1, n\}}$ are isomorphisms. ${ }^{5}$\\
The proof of Proposition 1.1.14 requires clever combinatorics, due to Joyal, and is deferred to Proposition D.4.6. Here, we enjoy its myriad consequences. Immediately:\\
1.1.15. Corollary. A quasi-category is a Kan complex if and only if its homotopy category is a groupoid.

Proof. If the homotopy category of a quasi-category is a groupoid, then all of its 1-simplices are isomorphisms, and Proposition 1.1.14 then implies that all inner and outer horns have fillers. Thus, the quasi-category is a Kan complex. Conversely, in a Kan complex, all outer horns can be filled and in particular fillers for the horns displayed in Definition 1.1.13 can be used to construct left and right inverses for any 1-simplex, which can be rectified to a single two-sided inverse by Lemma 1.1.12. \(\square\)

A quasi-category contains $A$ a canonical maximal sub Kan complex $A^{\simeq}$, the simplicial subset spanned by those 1-simplices that are isomorphisms. Just as the arrows in a quasi-category $A$ are represented by simplicial maps $2 \rightarrow A$ whose domain is the nerve of the free-living arrow, the isomorphisms in a quasi-category can be represented by diagrams $\mathbb{I} \rightarrow A$ whose domain, called the homotopy coherent isomorphism, is the nerve of the free-living isomorphism:

\footnotetext{${ }^{5}$ In the case $n=1$, no condition is needed on the horns; degenerate 1 -simplices define the required lifts.
}
1.1.16. Corollary. An arrow $f$ in a quasi-category $A$ is an isomorphism if and only if it extends to a homotopy coherent isomorphism\\
\includegraphics[max width=\textwidth, alt={}, center]{b375083c-8174-49a1-a017-e977744038d8-07_193_197_350_963}

Proof. If $f$ is an isomorphism, the map $f: 2 \rightarrow A$ lands in the maximal sub Kan complex contained in $A$ :\\
\includegraphics[max width=\textwidth, alt={}, center]{b375083c-8174-49a1-a017-e977744038d8-07_193_301_667_911}\\
By Exercise 1.1.v, the inclusion 2 ↪ II can be expressed as a sequential composite of pushouts of outer horn inclusions. Since $A^{\simeq}$ is a Kan complex, this shows that the required extension exists and in fact lands in $A^{\simeq} \subset A$. \(\square\)

The category of simplicial sets, like any category of presheaves, is cartesian closed. By the Yoneda lemma and the defining adjunction, an $n$-simplex in the exponential $Y^{X}$ corresponds to a simplicial map $X \times \Delta[n] \rightarrow Y$, and its faces and degeneracies are computed by precomposing in the simplex variable. Our next aim is to show that the quasi-categories define an exponential ideal in the simplicially enriched category of simplicial sets: if $X$ is a simplicial set and $A$ is a quasi-category, then $A^{X}$ is a quasi-category. We deduce this as a corollary of the "relative" version of this result involving certain maps called isofibrations that we now introduce.\\
1.1.17. Definition (isofibration). A simplicial map $f: A \rightarrow B$ between quasi-categories is an isofibration if it lifts against the inner horn inclusions, as displayed below-left, and also against the inclusion of either vertex into the free-living isomorphism II.\\
\includegraphics[max width=\textwidth, alt={}, center]{b375083c-8174-49a1-a017-e977744038d8-07_194_278_1587_717}\\
\includegraphics[max width=\textwidth, alt={}, center]{b375083c-8174-49a1-a017-e977744038d8-07_184_206_1594_1208}\\
To notationally distinguish the isofibrations, we depict them as arrows "→" with two heads.\\
Proposition 1.1.14 is subsumed by its relative analogue, proven as Theorem D.5.1:\\
1.1.18. Proposition (special outer horn lifting). Any isofibration between quasi-categories $f: A \rightarrow B$ admits lifts against those outer horns\\
\includegraphics[max width=\textwidth, alt={}, center]{b375083c-8174-49a1-a017-e977744038d8-07_214_279_2076_562}\\
\includegraphics[max width=\textwidth, alt={}, center]{b375083c-8174-49a1-a017-e977744038d8-07_214_520_2076_1054}\\
in which the edges $\left.g\right|_{\{0,1\}},\left.h\right|_{\{n-1, n\}},\left.k\right|_{\{0,1\}}$, and $\left.\ell\right|_{\{n-1, n\}}$ are isomorphisms.\\
1.1.19. Observation.

\begin{itemize}
  \item[(i)] For any simplicial set $X$, the unique map $X \rightarrow 1$ whose codomain is the terminal simplicial set is an isofibration if and only if $X$ is a quasi-category.
  \item[(ii)] Any collection of maps, such as the isofibrations, that is characterized by a right lifting property is automatically closed under composition, product, pullback, retract, and (inverse) limits of towers (see Lemma C.2.3).
  \item[(iii)] Combining (i) and (ii), if $A \rightarrow B$ is an isofibration, and $B$ is a quasi-category, then so is $A$.
  \item[(iv)] The isofibrations generalize the eponymous categorical notion. The nerve of any functor $f: A \rightarrow B$ between categories defines a map of simplicial sets that lifts against the inner horn inclusions. This map then defines an isofibration if and only if given any isomorphism in $B$ and specified object in $A$ lifting either its domain or codomain, there exists an isomorphism in $A$ with that domain or codomain lifting the isomorphism in $B$.
\end{itemize}

Much harder to establish is the stability of the isofibrations under the formation of "Leibniz ${ }^{6}$ exponentials" as displayed in (1.1.21). This is proven in Proposition D.5.2.\\
1.1.20. Proposition. If $i: X \hookrightarrow Y$ is a monomorphism and $f: A \rightarrow B$ is an isofibration, then the induced Leibniz exponential map $i \widehat{\hbar} f$\\
\includegraphics[max width=\textwidth, alt={}, center]{b375083c-8174-49a1-a017-e977744038d8-08_352_447_1038_842}

(1.1.21)\\
is again an isofibration. ${ }^{7}$\\
1.1.22. Corollary. If $X$ is a simplicial set and $A$ is a quasi-category, then $A^{X}$ is a quasi-category. Moreover, a 1-simplex in $A^{X}$ is an isomorphism if and only if its components at each vertex of $X$ are isomorphisms in $A$.

Proof. The first statement is a special case of Proposition 1.1.20 (see Exercise 1.1.vii), while the second statement is proven similarly by arguing with marked simplicial sets (see Corollary D.4.19). \(\square\)\\
1.1.23. Definition (equivalences of quasi-categories). A map $f: A \rightarrow B$ between quasi-categories is an equivalence if it extends to the data of a "homotopy equivalence" with the free-living isomorphism II serving as the interval: that is, if there exist maps $g: B \rightarrow A$,\\
\includegraphics[max width=\textwidth, alt={}, center]{b375083c-8174-49a1-a017-e977744038d8-08_315_245_1878_680}\\
and\\
\includegraphics[max width=\textwidth, alt={}, center]{b375083c-8174-49a1-a017-e977744038d8-08_311_235_1879_1214}\\
We write "→" to decorate equivalences and $A \simeq B$ to indicate the presence of an equivalence $A \xrightarrow{3}$.

\footnotetext{${ }^{6}$ The name alludes to the Leibniz rule in differential calculus, or more specifically to the identification of the domain of the Leibniz product of Lemma D.3.1 with the boundary of the prism (see Definition C.2.8 and Remark D.3.2).\\
${ }^{7}$ Degenerate cases of this result, taking $X=\varnothing$ or $B=1$, imply that the other six maps in this diagram are also isofibrations (see Exercise 1.1.vii).
}
1.1.24. Remark. If $f: A \rightarrow B$ is an equivalence of quasi-categories, then the functor $\mathrm{h} f: \mathrm{h} A \rightarrow \mathrm{~h} B$ is an equivalence of categories, where the data displayed above defines an equivalence inverse $\mathrm{h} g: \mathrm{h} B \rightarrow$ $\mathrm{h} A$ and natural isomorphisms encoded by the composite ${ }^{8}$ functors
$$
\mathrm{h} A \xrightarrow{\mathrm{~h} \alpha} \mathrm{~h}\left(A^{\mathbb{I}}\right) \longrightarrow(\mathrm{h} A)^{\mathbb{I}} \quad \mathrm{h} B \xrightarrow{\mathrm{~h} \beta} \mathrm{~h}\left(B^{\mathbb{I}}\right) \longrightarrow(\mathrm{h} B)^{\mathbb{I}}
$$
1.1.25. Definition. A map $f: X \rightarrow Y$ between simplicial sets is a trivial fibration if it admits lifts against the boundary inclusions for all simplices\\
\includegraphics[max width=\textwidth, alt={}, center]{b375083c-8174-49a1-a017-e977744038d8-09_193_540_638_799}

(1.1.26)\\
We write " ${ }^{2}$ »" to decorate trivial fibrations.\\
1.1.27. Remark. The simplex boundary inclusions $\partial \Delta[n] \hookrightarrow \Delta[n]$ "cellularly generate" the monomorphisms of simplicial sets (see Definition C.2.4 and Lemma C.5.9). Hence the dual of Lemma C.2.3 implies that trivial fibrations lift against any monomorphism between simplicial sets. In particular, it follows that any trivial fibration $X \xrightarrow{\curvearrowright} Y$ is a split epimorphism.

The notation " $\leadsto$ " is suggestive: the trivial fibrations between quasi-categories are exactly those maps that are both isofibrations and equivalences. This can be proven by a relatively standard although rather technical argument in simplicial homotopy theory, appearing as Proposition D.5.6.\\
1.1.28. Proposition. For a map $f: A \rightarrow B$ between quasi-categories the following are equivalent:

\begin{itemize}
  \item[(i)] $f$ is a trivial fibration
  \item[(ii)] $f$ is both an isofibration and an equivalence
  \item[(iii)] $f$ is a split fiber homotopy equivalence: an isofibration admitting a section $s$ that is also an equivalence inverse via a homotopy $\alpha$ from $\mathrm{id}_{A}$ to $s f$ that composes with $f$ to the constant homotopy from $f$ to $f$.\\
\includegraphics[max width=\textwidth, alt={}, center]{b375083c-8174-49a1-a017-e977744038d8-09_208_455_1575_839}
\end{itemize}

As a class characterized by a right lifting property, the trivial fibrations are also closed under composition, product, pullback, limits of towers, and contain the isomorphisms. The stability of these maps under Leibniz exponentiation is proven along with Proposition 1.1.20 in Proposition D.5.2.\\
1.1.29. Proposition. If $i: X \rightarrow Y$ is a monomorphism and $f: A \rightarrow B$ is an isofibration, then if either $f$ is a trivial fibration or if $i$ is in the class cellularly generated by the inner horn inclusions and the map $\mathbb{1} \leftrightarrow \mathbb{I}$ then the induced Leibniz exponential map

$$
A^{Y} \xrightarrow{i \hbar f} B^{Y} \times_{B^{X}} A^{X}
$$

a trivial fibration.

\footnotetext{${ }^{8}$ Note that $\mathrm{h}\left(A^{\mathbb{I}}\right) \neq(\mathrm{h} A)^{\mathbb{I}}$ in general. Objects in the latter are homotopy classes of isomorphisms in $A$, while objects in the former are homotopy coherent isomorphisms, given by a specified 1-simplex in $A$, a specified inverse 1-simplex, together with an infinite tower of coherence data indexed by the nondegenerate simplices in $\mathbb{I}$.
}To illustrate the utility of these Leibniz stability results, we give an "internal" or "synthetic" characterization of the Kan complexes.\\
1.1.30. Lemma. A quasi-category $A$ is a Kan complex if and only if the map $A^{\mathbb{I}} \rightarrow A^{2}$ induced by the inclusion $2 \hookrightarrow \mathbb{I}$ is a trivial fibration.

Note that Proposition 1.1.20 implies that $A^{\mathbb{I}} \rightarrow A^{2}$ is an isofibration.\\
Proof. The lifting property that characterizes trivial fibrations transposes to another lifting property, displayed below-right\\
\includegraphics[max width=\textwidth, alt={}, center]{b375083c-8174-49a1-a017-e977744038d8-10_219_1186_702_469}\\
that asserts that $A$ admits extensions along maps formed by taking the Leibniz product - also known as the pushout product - of a simplex boundary inclusion $\partial \Delta[n] \hookrightarrow \Delta[n]$ with the inclusion $2 \hookrightarrow \mathbb{I}$. By Exercise 1.1.v(ii) the inclusion $2 \hookrightarrow \mathbb{I}$ is a sequential composite of pushouts of left outer horn inclusions. By Corollary D.3.11, a key step along the way to the proofs of Propositions 1.1.20 and 1.1.29, it follows that the Leibniz product is also a sequential composite of pushouts of left and inner horn inclusions. If $A$ is a Kan complex, then the extensions displayed above right exist, and, by transposing, the map $A^{\mathbb{I}} \rightarrow A^{2}$ is a trivial fibration.

Conversely, if $A^{\mathbb{I}} \xrightarrow{{ }^{2}}$ is a trivial fibration then in particular it is surjective on vertices. Thus every arrow in $A$ is an isomorphism, and Corollary 1.1.15 tells us that $A$ must be a Kan complex. \(\square\)\\
1.1.31. Digression (the Joyal model structure). The category of simplicial sets bears a Quillen model structure, in the sense of Definition C.3.1, whose fibrant objects are exactly the quasi-categories and in which all objects are cofibrant. Between fibrant objects, the fibrations, weak equivalences, and trivial fibrations are precisely the isofibrations, equivalences, and trivial fibrations just introduced. Proposition 1.1.28 proves that the trivial fibrations are the intersection of the fibrations and the weak equivalences. Propositions 1.1.20 and 1.1.29 reflect the fact that the Joyal model structure is a cartesian closed model category, satisfying the additional axioms of Definition C.3.10.

We decline to elaborate further on the Joyal model structure for quasi-categories since we have highlighted all of the features that we need. The results enumerated here suffice to show that the category of quasi-categories defines an $\infty$-cosmos, a concept to which we now turn.

\section*{Exercises.}
1.1.i. Exercise ([103, §1.5]). Given any cosimplicial object $C: \boldsymbol{\Delta} \rightarrow \mathcal{E}$ valued in any category $\mathcal{E}$, there is an associated nerve functor $N_{C}$ defined by:\\
\includegraphics[max width=\textwidth, alt={}, center]{b375083c-8174-49a1-a017-e977744038d8-10_149_360_2148_625}\\
\includegraphics[max width=\textwidth, alt={}, center]{b375083c-8174-49a1-a017-e977744038d8-10_229_406_2144_1091}

By construction $n$-simplices in $N_{C} E$ correspond to maps $C^{n} \rightarrow E$ in $\mathcal{E}$. Show that if $\mathcal{E}$ is cocomplete, then $N_{C}$ has a left adjoint defined as the left Kan extension of the functor $C$ along the Yoneda embedding よ: $\Delta \hookrightarrow$ sSet. This gives a second proof of Proposition 1.1.11.\\[0pt]
1.1.ii. Exercise (Boardman-Vogt [21]). Consider the set of 1-simplices in a quasi-category with initial vertex $a$ and final vertex $b$.

\begin{itemize}
  \item[(i)] Prove that the relation defined by $f \sim g$ if and only if there exists a 2 -simplex with boundary\\
\includegraphics[max width=\textwidth, alt={}, center]{b375083c-8174-49a1-a017-e977744038d8-11_157_233_586_385}\\
is an equivalence relation.
  \item[(ii)] Prove that the relation defined by $f \sim g$ if and only if there exists a 2-simplex with boundary\\
\includegraphics[max width=\textwidth, alt={}, center]{b375083c-8174-49a1-a017-e977744038d8-11_147_233_804_386}\\
is an equivalence relation.
  \item[(iii)] Prove that the equivalence relations defined by (i) and (ii) are the same.
\end{itemize}

This proves Lemma 1.1.9.\\
1.1.iii. Exercise. Consider the free category on the reflexive directed graph

$$
A_{1} \xrightarrow[\cdot \delta^{0}]{\xrightarrow{{ }^{+} \cdot \delta^{1}}-} A_{0},
$$

underlying a quasi-category $A$.

\begin{itemize}
  \item[(i)] Consider the binary relation that identifies sequences of composable 1-simplices with common source and common target whenever there exists a simplex of $A$ in which the sequences of 1-simplices define two paths from its initial vertex to its final vertex. Prove that this relation is stable under pre- and post-composition with 1-simplices and conclude that its transitive closure is a congruence: an equivalence relation that is closed under pre- and post-composition. ${ }^{9}$
  \item[(ii)] Consider the congruence relation generated by imposing a composition relation $h=g \circ f$ witnessed by 2 -simplices\\
\includegraphics[max width=\textwidth, alt={}, center]{b375083c-8174-49a1-a017-e977744038d8-11_145_283_1727_919}\\
and prove that this coincides with the relation considered in (i).
  \item[(iii)] In the congruence relations of (i) and (ii), prove that every sequence of composable 1-simplices in $A$ is equivalent to a single 1-simplex. Conclude that every morphism in the quotient of the free category by this congruence relation is represented by a 1 -simplex in $A$.
  \item[(iv)] Prove that for any triple of 1-simplices $f, g, h$ in $A, h=g \circ f$ in the homotopy category $\mathrm{h} A$ of Definition 1.1.10 if and only if there exists a 2 -simplex with boundary\\
\includegraphics[max width=\textwidth, alt={}, center]{b375083c-8174-49a1-a017-e977744038d8-11_148_282_2198_920}
\end{itemize}

\footnotetext{${ }^{9}$ Given a congruence relation on the hom-sets of a 1-category, the quotient category can be formed by quotienting each hom-set (see [81, §II.8]).
}This proves Lemma 1.1.12.\\
1.1.iv. Exercise. Show that any quasi-category in which inner horns admit unique fillers is isomorphic to the nerve of its homotopy category.\\
1.1. v. Exercise. Let II be the nerve of the free-living isomorphism.

\begin{itemize}
  \item[(i)] Prove that II contains exactly two nondegenerate simplices in each dimension.
  \item[(ii)] Inductively build II from 2 by expressing the inclusion $2 \hookrightarrow \mathbb{I}$ as a sequential composite of pushouts of left outer horn inclusions ${ }^{10} \Lambda^{0}[n] \hookrightarrow \Delta[n]$, one in each dimension starting with $n=2 .{ }^{11}$
\end{itemize}

\href{http://1.1.vi}{1.1.vi}. Exercise. Prove the relative version of Corollary 1.1.16: for any isofibration $p: A \rightarrow B$ between quasi-categories and any $f: 2 \rightarrow A$ that defines an isomorphism in $A$ any homotopy coherent isomorphism in $B$ extending $p f$ lifts to a homotopy coherent isomorphism in $A$ extending $f$.\\
\includegraphics[max width=\textwidth, alt={}, center]{b375083c-8174-49a1-a017-e977744038d8-12_195_206_888_962}\\
1.1.vii. Exercise. Specialize Proposition 1.1.20 to prove the following:

\begin{itemize}
  \item[(i)] If $A$ is a quasi-category and $X$ is a simplicial set then $A^{X}$ is a quasi-category.
  \item[(ii)] If $A$ is a quasi-category and $X \hookrightarrow Y$ is a monomorphism then $A^{Y} \rightarrow A^{X}$ is an isofibration.
  \item[(iii)] If $A \rightarrow B$ is an isofibration and $X$ is a simplicial set then $A^{X} \rightarrow B^{X}$ is an isofibration.
\end{itemize}

1.1.viii. Exercise. Anticipating Lemma 1.2.17:

\begin{itemize}
  \item[(i)] Prove that the equivalences defined in Definition 1.1.23 are closed under retracts.
  \item[(ii)] Prove that the equivalences defined in Definition 1.1.23 satisfy the 2-of-3 property.
\end{itemize}

1.1.ix. Exercise. Prove that if $f: X \leadsto Y$ is a trivial fibration between quasi-categories then the functor $\mathrm{h} f: \mathrm{h} X \leadsto \mathrm{~h} Y$ is a surjective equivalence of categories.

\section*{1.2. $\infty$-Cosmoi}
In §1.1, we presented "analytic" proofs of a few of the basic facts about quasi-categories. The category theory of quasi-categories can be developed in a similar style, but we aim instead to develop the "synthetic" theory of infinite-dimensional categories, so that our results apply to many models at once. To achieve this, our strategy is not to axiomatize what infinite-dimensional categories are, but rather to axiomatize the categorical "universe" in which they live.

The definition of an $\infty$-cosmos abstracts the properties of the category of quasi-categories together with the isofibrations, equivalences, and trivial fibrations introduced in §1.1. ${ }^{12}$ First, the category of quasi-categories is enriched over the category of simplicial sets - the set of morphisms from $A$ to $B$ coincides with the set of vertices of the simplicial set $B^{A}$ - and moreover these hom spaces are all

\footnotetext{${ }^{10}$ By the duality described in Definition 1.2.25, the right outer horn inclusions $\Lambda^{n}[n] \hookrightarrow \Delta[n]$ can be used instead.\\
${ }^{11}$ This decomposition of the inclusion $2 \hookrightarrow \mathbb{I}$ reveals which data extends homotopically uniquely to a homotopy coherent isomorphism. For instance, the 1- and 2-simplices of Definition 1.1.13 together with a single 3-simplex that has these as its outer faces with its inner faces degenerate. Homotopy type theorists refer to this data as a half adjoint equivalence [125, §4.2].\\
${ }^{12}$ Metaphorical allusions aside, our $\infty$-cosmoi resemble the fibrational cosmoi of Street [117].
}
quasi-categories. Second, certain limit constructions that can be defined in the underlying unenriched category of quasi-categories satisfy universal properties relative to this simplicial enrichment, with the usual isomorphism of sets extending to an isomorphism of simplicial sets. And finally, the isofibrations, equivalences, and trivial fibrations satisfy properties that are familiar from abstract homotopy theory, forming a category of fibrant objects à la Brown [23] (see §C.1). In particular, the use of isofibrations in diagrams guarantees that their strict limits are equivalence invariant, so we can take advantage of up-to-isomorphism universal properties and strict functoriality of these constructions while still working "homotopically."

As explained in Digression 1.2.13, there are a variety of models of infinite-dimensional categories for which the category of " $\infty$-categories," as we call them, and " $\infty$-functors" between them is enriched over quasi-categories and admits classes of isofibrations, equivalences, and trivial fibrations satisfying analogous properties. This motivates the following axiomatization:\\
1.2.1. Definition $(\infty-\operatorname{cosmos})$. An $\infty$-cosmos $\mathcal{K}$ is a category that is enriched over quasi-categories, ${ }^{13}$ meaning in particular that

\begin{itemize}
  \item[-] its morphisms $f: A \rightarrow B$ define the vertices of a quasi-category denoted $\operatorname{Fun}(A, B)$ and referred to as a functor space,
\end{itemize}

that is also equipped with a specified collection of maps that we call isofibrations and denote by "→" satisfying the following two axioms:

\begin{itemize}
  \item[]
  \begin{itemize}
    \item[(i)] (completeness) The quasi-categorically enriched category $\mathcal{K}$ possesses a terminal object, small products, pullbacks of isofibrations, limits of countable towers of isofibrations, and cotensors with simplicial sets, each of these limit notions satisfying a universal property that is enriched over simplicial sets. ${ }^{14}$
    \item[(ii)] (isofibrations) The isofibrations contain all isomorphisms and any map whose codomain is the terminal object; are closed under composition, product, pullback, forming inverse limits of towers, and Leibniz cotensors with monomorphisms of simplicial sets; and have the property that if $f: A \rightarrow B$ is an isofibration and $X$ is any object then $\operatorname{Fun}(X, A) \rightarrow \operatorname{Fun}(X, B)$ is an isofibration of quasi-categories.
  \end{itemize}
\end{itemize}

For ease of reference, we refer to the simplicially enriched limits of diagrams of isofibrations enumerated in (i) as the cosmological limit notions.\\
1.2.2. Definition. In an $\infty$-cosmos $\mathcal{K}$, a morphism $f: A \rightarrow B$ is

\begin{itemize}
  \item[-] an equivalence just when the induced map $f_{*}: \operatorname{Fun}(X, A) \xrightarrow{\sim \operatorname{Fun}(X, B) \text { on functor spaces is an }}$ equivalence of quasi-categories for all $X \in \mathcal{K}$, and
  \item[-] a trivial fibration just when $f$ is both an isofibration and an equivalence.
\end{itemize}

These classes are denoted by "→" and "→»", respectively.\\
Put more concisely, one might say that an $\infty$-cosmos is a "quasi-categorically enriched category of fibrant objects" (see Definition C.1.1 and Example C.1.3).\\
1.2.3. Convention ( $\infty$-category, as a technical term). Henceforth, we recast $\infty$-category as a technical term to refer to an object in an arbitrary ambient $\infty$-cosmos. Similarly, we use the term $\infty$-functor - or more commonly the elision "functor" - to refer to a morphism $f: A \rightarrow B$ in an $\infty$-cosmos. This

\footnotetext{${ }^{13}$ This is to say $\mathcal{K}$ is a simplicially enriched category (see Digression 1.2.4) whose hom spaces are all quasi-categories.\\
${ }^{14}$ We elaborate on these simplicially enriched limits in Digression 1.2.6.
}
explains why we refer to the quasi-category $\operatorname{Fun}(A, B)$ between two $\infty$-categories in an $\infty$-cosmos as a "functor space": its vertices are the $(\infty-)$ functors from $A$ to $B$.\\
1.2.4. Digression (simplicial categories, § A.2). A simplicial category $\mathcal{A}$ is given by categories $\mathcal{A}_{n}$, with a common set of objects and whose arrows are called $n$-arrows, that assemble into a diagram $\boldsymbol{\Delta}^{\mathrm{op}} \rightarrow$ Cat of identity-on-objects functors

\begin{equation*}
\cdots \mathcal{A}_{3} \underset{-\cdot \delta^{1}}{\stackrel{-\cdot \delta^{3}}{\rightleftarrows} \cdot \delta^{2} \longrightarrow} \mathcal{A}_{2} \underset{-\cdot \delta^{1}}{\stackrel{-\delta^{2}}{\leftrightarrows} \cdot \delta^{1} \longrightarrow} \mathcal{A}_{1} \underset{-\cdot \delta^{0}}{\stackrel{-\cdot \delta^{1}}{\leftrightarrows}} \underset{\rightarrow}{\stackrel{\leftarrow \cdot \delta^{0}}{\longrightarrow}} \mathcal{A}_{0}=: \mathcal{A} \tag{1.2.5}
\end{equation*}

The category $\mathcal{A}_{0}$ of 0-arrows is the underlying category of the simplicial category $\mathcal{A}$, which forgets the higher dimensional simplicial structure.

The data of a simplicial category can equivalently be encoded by a simplicially enriched category with a set of objects and a simplicial set $\mathcal{A}(x, y)$ of morphisms between each ordered pair of objects: an $n$-arrow in $\mathcal{A}_{n}$ from $x$ to $y$ corresponds to an $n$-simplex in $\mathcal{A}(x, y)$ (see Exercise 1.2.i). Each endo-hom space contains a distinguished identity 0-arrow (the degenerate images of which define the corresponding identity $n$-arrows) and composition is required to define a simplicial map

$$
\mathcal{A}(y, z) \times \mathcal{A}(x, y) \longrightarrow \mathcal{A}(x, z)
$$

the single map encoding the compositions in each of the categories $\mathcal{A}_{n}$ and also the functoriality of the diagram (1.2.5). The composition is required to be associative and unital, in a sense expressed by the commutative diagrams of simplicial sets

$$
\begin{gathered}
\mathcal{A}(y, z) \times \mathcal{A}(x, y) \times \mathcal{A}(w, x) \xrightarrow{\circ \times \mathrm{id}} \mathcal{A}(x, z) \times \mathcal{A}(w, x) \\
\downarrow \\
\begin{array}{l}
\mathrm{id} \times \circ \downarrow \\
\mathcal{A}(y, z) \times \mathcal{A}(w, y) \longrightarrow
\end{array}
\end{gathered}
$$

\includegraphics[max width=\textwidth, alt={}, center]{b375083c-8174-49a1-a017-e977744038d8-14_225_732_1393_1144}\\
On account of the equivalence between these two presentations, the terms "simplicial category" and "simplicially enriched category" are generally taken to be synonyms. ${ }^{15}$

In particular, the underlying category $\mathcal{K}_{0}$ of an $\infty$-cosmos $\mathcal{K}$ is the category whose objects are the $\infty$-categories in $\mathcal{K}$ and whose morphisms are the 0-arrows, i.e., the vertices in the functor spaces. In all of the examples to appear in what follows, this recovers the expected category of $\boldsymbol{\infty}$-categories in a particular model and functors between them.\\
1.2.6. Digression (simplicially enriched limits, §A.4-A.5). Let $\mathcal{A}$ be a simplicial category. The cotensor of an object $A \in \mathcal{A}$ by a simplicial set $U$ is characterized by a natural isomorphism of simplicial sets


\begin{equation*}
\mathcal{A}\left(X, A^{U}\right) \cong \mathcal{A}(X, A)^{U} \tag{1.2.7}
\end{equation*}


\footnotetext{${ }^{15}$ The phrase "simplicial object in Cat" is reserved for the more general yet less common notion of a diagram $\boldsymbol{\Delta}^{\mathrm{op}} \rightarrow$ Cat that is not necessarily comprised of identity-on-objects functors.
}Assuming such objects exist, the simplicial cotensor defines a bifunctor

$$
\begin{aligned}
s \mathcal{S} e t^{\text {op }} \times \mathcal{A} & \longrightarrow \mathcal{A} \\
(U, A) & \longmapsto A^{U}
\end{aligned}
$$

in a unique way making the isomorphism (1.2.7) natural in $U$ and $A$ as well.\\
The other simplicial limit notions postulated by axiom 1.2.1(i) are conical, which is the term used for ordinary 1-categorical limit shapes that satisfy an enriched analog of the usual universal property (see Definition A.5.2). Such limits also define limits in the underlying category, but the usual universal property is strengthened. By applying the covariant representable functor $\mathcal{A}(X,-): \mathcal{A}_{0} \rightarrow s \mathcal{S} e t$ to a limit cone $\left(\lim _{j \in J} A_{j} \rightarrow A_{j}\right)_{j \in J}$ in $\mathcal{A}_{0}$, we obtain a natural comparison map


\begin{equation*}
\mathcal{A}\left(X, \lim _{j \in J} A_{j}\right) \longrightarrow \lim _{j \in J} \mathcal{A}\left(X, A_{j}\right) . \tag{1.2.8}
\end{equation*}


We say that $\lim _{j \in J} A_{j}$ defines a simplicially enriched limit if and only if (1.2.8) is an isomorphism of simplicial sets for all $X \in \mathcal{A}$.

The general theory of enriched categories is reviewed in Appendix A.\\
1.2.9. Preview (flexible weighted limits in $\infty$-cosmoi). The axiom 1.2.1(i) implies that any $\infty$-cosmos $\mathcal{K}$ admits all flexible limits, a much larger class of simplicially enriched "weighted" limits (see Definition 6.2.1 and Proposition 6.2.8).

We quickly introduce the three examples of $\infty$-cosmoi that are most easily absorbed, deferring more sophisticated examples to the end of this section. The first of these is the prototypical $\infty$-cosmos.\\
1.2.10. Proposition (the $\infty$-cosmos of quasi-categories). The full subcategory QCat $\subset$ sSet of quasi-cat-egories defines an $\infty$-cosmos in which the isofibrations, equivalences, and trivial fibrations coincide with the classes already bearing these names.

Proof. The subcategory QCat $\subset$ sSet inherits its simplicial enrichment from the cartesian closed category of simplicial sets: by Proposition 1.1.20, whenever $A$ and $B$ are quasi-categories, $\operatorname{Fun}(A, B):=B^{A}$ is again a quasi-category.

The cosmological limits postulated in 1.2.1(i) exist in the ambient category of simplicial sets. ${ }^{16}$ For instance, the defining universal property of the simplicial cotensor (1.2.7) is satisfied by the exponentials of simplicial sets. Moreover, since the category of simplicial sets is cartesian closed, each of the conical limits is simplicially enriched in the sense discussed in Digression 1.2.6 (see Exercise 1.2.ii and Proposition A.5.4).

We now argue that the full subcategory of quasi-categories inherits all these limit notions and at the same time establish the stability of the isofibrations under the formation of these limits. In fact, this latter property helps to prove the former. To see this, note that a simplicial set is a quasi-category if and only if the map from it to the point is an isofibration. More generally, if the codomain of any isofibration is a quasi-category then its domain must be as well. So if any of the maps in a limit cone over a diagram of quasi-categories are isofibrations, then it follows that the limit is itself a quasi-category.

Since the isofibrations are characterized by a right lifting property, Lemma C.2.3 implies that the isofibrations contains all isomorphism and are closed under composition, product, pullback, and forming inverse limits of towers. In particular, the full subcategory of quasi-categories possesses these

\footnotetext{${ }^{16}$ Any category of presheaves is cartesian closed, complete, and cocomplete - a "cosmos" in the sense of Bénabou.
}
limits. This verifies all of the axioms of 1.2.1(i) and 1.2.1(ii) except for the last two: Leibniz closure and closure under exponentiation $(-)^{X}$. These last closure properties are established in Proposition 1.1.20, and in fact by Exercise 1.1.vii, the former subsumes the latter . This completes the verification of the $\infty$-cosmos axioms.

It remains to check that the equivalences and trivial fibrations coincide with those maps defined by 1.1.23 and 1.1.25. By Proposition 1.1.28 the latter coincidence follows from the former, so it remains only to show that the equivalences of 1.1.23 coincide with the representably defined equivalences: those maps of quasi-categories $f: A \rightarrow B$ for which $A^{X} \rightarrow B^{X}$ is an equivalence of quasi-categories in the sense of Definition 1.1.23. Taking $X=\Delta[0]$, we see immediately that representably defined equivalences are equivalences, and the converse holds since the exponential $(-)^{X}$ preserves the data defining a simplicial homotopy. \(\square\)

Two further examples fit into a common paradigm: both arise as full subcategories of the $\infty$-cosmos of quasi-categories and inherit their $\infty$-cosmos structures from this inclusion (see Lemma 6.1.4). But it is also instructive, and ultimately takes less work, to describe the resulting $\infty$-cosmos structures directly.

\begin{figure}[h]
\begin{center}
\captionsetup{labelformat=empty}
\caption{1.2.11. Proposition (the $\infty$-cosmos of categories). The category Cat of 1 -categories defines an $\infty$-cosmos whose isofibrations are the isofibrations: functors satisfying the displayed right lifting property:}
  \includegraphics[alt={},max width=\textwidth]{b375083c-8174-49a1-a017-e977744038d8-16_172_207_1156_962}
\end{center}
\end{figure}

The equivalences are the equivalences of categories and the trivial fibrations are surjective equivalences: equivalences of categories that are also surjective on objects.

Proof. It is well-known that the 2-category of categories is complete (and in fact also cocomplete) as a Cat-enriched category (see Definition A.6.17 or [67]). The categorically enriched category of categories becomes a quasi-categorically enriched category by applying the nerve functor to the homcategories (see §A.7). Since the nerve functor is a right adjoint, it follows formally that these 2-categorical limits become simplicially enriched limits. In particular, as proscribed in Proposition A.7.8, the cotensor of a category $A$ by a simplicial set $U$ is defined to be the functor category $A^{\text {hU }}$. This completes the verification of axiom (i).

Since the class of isofibrations is characterized by a right lifting property, Lemma C.2.3 implies that the isofibrations are closed under all of the limit constructions of 1.2.1(ii) except for the last two, and by Exercise 1.1.vii, the Leibniz closure subsumes the closure under exponentiation.

To verify that isofibrations of categories $f: A \rightarrow B$ are stable under forming Leibniz cotensors with monomorphisms of simplicial sets $i: U \hookrightarrow V$, we must solve the lifting problem below-left\\
\includegraphics[max width=\textwidth, alt={}, center]{b375083c-8174-49a1-a017-e977744038d8-16_234_975_2081_575}\\
which transposes to the lifting problem above-right, which we can solve by hand. Here the map $\beta$ defines a natural isomorphism between $f s: \mathrm{h} V \rightarrow B$ and a second functor. Our task is to lift this\\
to a natural isomorphism $\gamma$ from $s$ to another functor that extends the natural isomorphism $\alpha$ along $\mathrm{h} i: \mathrm{h} U \rightarrow \mathrm{~h} V$. Note this functor hi need not be an inclusion, but it is injective on objects, which is enough.

We define the components of $\gamma$ by cases. If an object $v \in \mathrm{~h} V$ is equal to $i(u)$ for some $u \in \mathrm{~h} U$ define $\gamma_{i(u)}:=\alpha_{u}$; otherwise, use the fact that $f$ is an isofibration to define $\gamma_{v}$ to be any lift of the isomorphism $\beta_{v}$ to an isomorphism in $A$ with domain $s(v)$. The data of the map $\gamma: \mathrm{h} V \times \mathbb{I} \rightarrow A$ also entails the specification of the functor $\mathrm{h} V \rightarrow A$ that is the codomain of the natural isomorphism $\gamma$. On objects, this functor is given by $v \mapsto \operatorname{cod}\left(\gamma_{v}\right)$. On morphisms, this functor defined in the unique way that makes $\gamma$ into a natural transformation:

$$
\left(k: v \rightarrow v^{\prime}\right) \mapsto \gamma_{v^{\prime}} \circ s(k) \circ \gamma_{v}^{-1} .
$$

This completes the proof that $C a t$ defines an $\infty$-cosmos. Since the nerve of a functor category, such as $A^{\mathbb{I}}$, is isomorphic to the exponential between their nerves, the equivalences of categories coincide with the equivalences of Definition 1.1.23. It follows that the equivalences in the $\infty$-cosmos of categories coincide with equivalences of categories, and since the surjective equivalences are the intersection of the equivalences and the isofibrations, this completes the proof. \(\square\)\\
1.2.12. Proposition (the $\infty$-cosmos of Kan complexes). The category Kan of Kan complexes defines an $\infty$-cosmos whose isofibrations are the Kan fibrations: maps that lift against all horn inclusions $\Lambda^{k}[n] \hookrightarrow \Delta[n]$ for $n \geq 1$ and $0 \leq k \leq n$.

The proof proceeds along the lines of Lemma 6.1.4. We show that the subcategory of Kan complexes inherits an $\infty$-cosmos structure by restricting structure from the $\infty$-cosmos of quasi-categories.

Proof. By Proposition 1.1.18, an isofibration between Kan complexes is a Kan fibration. Conversely, since the homotopy coherent isomorphism $\mathbb{I}$ can be built from the point $\mathbb{1}$ by attaching fillers to a sequence of outer horns, all Kan fibrations define isofibrations. This shows that between Kan complexes, isofibrations and Kan fibrations coincide. So to show that the category of Kan complexes inherits an $\infty$-cosmos structure by restriction from the $\infty$-cosmos of quasi-categories, we need only verify that the full subcategory Kan ↪ QCat is closed under all of the limit constructions of axiom 1.2.1(i). For the conical limits, the argument mirrors the one given in the proof of Proposition 1.2.10, while the closure under cotensors is a consequence of Corollary D.3.11, which implies that the Kan complexes also define an exponential ideal in the category of simplicial sets. The remaining axiom 1.2.1(ii) is inherited from the analogous properties established for quasi-categories in Proposition 1.2.10. \(\square\)

We mention a common source of $\infty$-cosmoi found in nature to build intuition for readers familiar with Quillen's model categories, a popular framework for abstract homotopy theory, but reassure newcomers that model categories are not needed outside of Appendix E where these results are proven.\\
1.2.13. Digression (a source of $\infty$-cosmoi in nature). As explained in §E.1, certain easily described properties of a model category imply that the full subcategory of fibrant objects defines an $\infty$-cosmos whose isofibrations, equivalences, and trivial fibrations are the fibrations, weak equivalences, and trivial fibrations between fibrant objects. Namely, any model category that is enriched as such over the Joyal model structure on simplicial sets in which all fibrant objects are cofibrant presents an $\infty$-cosmos (see Proposition E.1.1). This model-categorical enrichment over quasi-categories can be defined when the model category is cartesian closed and equipped with a right Quillen adjoint to the Joyal model structure on simplicial sets whose left adjoint preserves finite products (see Corollary E.1.4). In this case, the right adjoint becomes the underlying quasi-category functor (see Proposition 1.3.4(ii)) and\\
the $\infty$-cosmoi so-produced is cartesian closed (see Definition 1.2.23). The $\infty$-cosmoi listed in Example 1.2.24 all arise in this way.

The following results are consequences of the axioms of Definition 1.2.1. To begin, observe that the trivial fibrations enjoy the same stability properties satisfied by the isofibrations.\\
1.2.14. Lemma (stability of trivial fibrations). The trivial fibrations in an $\infty$-cosmos define a subcategory containing the isomorphisms and are stable under product, pullback, and forming inverse limits of towers. Moreover, the Leibniz cotensors of any trivial fibration with a monomorphism of simplicial sets is a trivial fibration as is the Leibniz cotensor of an isofibration with a map in the class cellularly generated by the inner horn inclusions and the map $\mathbb{1} \hookrightarrow \mathbb{I}$, and if $E \leadsto B$ is a trivial fibration then so is $\operatorname{Fun}(X, E) \leadsto \operatorname{Fun}(X, B)$.

Proof. We prove these statements in the reverse order. By axiom 1.2.1(ii) and the definition of the trivial fibrations in an $\infty$-cosmos, we know that if $E \xrightarrow{ } B$ is a trivial fibration then $\operatorname{Fun}(X, E) \xrightarrow{\leadsto}$ $\operatorname{Fun}(X, B)$ is both an isofibration and an equivalence, and hence by Proposition 1.1.28 a trivial fibration. For stability under the remaining constructions, we know in each case that the maps in question are isofibrations in the $\infty$-cosmos; it remains to show only that the maps are also equivalences. The equivalences in an $\infty$-cosmos are defined to be the maps that $\operatorname{Fun}(X,-)$ carries to equivalences of quasi-categories, so it suffices to verify that trivial fibrations of quasi-categories satisfy the corresponding stability properties. For the Leibniz stability properties, this is established in Proposition 1.1.29, while the remaining properties are covered by Lemma C.2.3. \(\square\)

By a Yoneda-style argument, the "homotopy equivalence" characterization of the equivalences in the $\infty$-cosmos of quasi-categories of Definition 1.1.23 extends to an analogous characterization of the equivalences in any $\infty$-cosmos:\\
1.2.15. Lemma (equivalences are homotopy equivalences). A map $f: A \rightarrow B$ between $\infty$-categories in an $\infty$-cosmos $\mathcal{K}$ is an equivalence if and only if it extends to the data of a "homotopy equivalence" with the free-living isomorphism II serving as the interval: that is, if there exist maps $g: B \rightarrow A$\\
\includegraphics[max width=\textwidth, alt={}, center]{b375083c-8174-49a1-a017-e977744038d8-18_317_767_1550_682}

(1.2.16)\\
in the $\infty$-cosmos.\\
Proof. By hypothesis, if $f: A \rightarrow B$ defines an equivalence in the $\infty$-cosmos $\mathcal{K}$ then the induced map on post-composition $f_{*}: \operatorname{Fun}(B, A) \xrightarrow{\sim \operatorname{Fun}(B, B) \text { is an equivalence of quasi-categories in the sense }}$ of Definition 1.1.23. Evaluating the inverse equivalence $\widetilde{g}: \operatorname{Fun}(B, B) \xrightarrow{\Rightarrow} \operatorname{Fun}(B, A)$ and homotopy $\widetilde{\beta}: \operatorname{Fun}(B, B) \rightarrow \operatorname{Fun}(B, B)^{\mathbb{I}}$ at the 0-arrow $\operatorname{id}_{B} \in \operatorname{Fun}(B, B)$, we obtain a 0-arrow $g: B \rightarrow A$ together with an isomorphism $\beta: \mathbb{I} \rightarrow \operatorname{Fun}(B, B)$ from the composite $f g$ to $\operatorname{id}_{B}$. By the defining universal property of the cotensor (1.2.7), this isomorphism internalizes to define the map $\beta: B \rightarrow B^{\mathbb{I}}$ in $\mathcal{K}$ displayed on the right of (1.2.16).

Now the hypothesis that $f$ is an equivalence also provides an equivalence of quasi-categories $f_{*}: \operatorname{Fun}(A, A) \xrightarrow{\sim} \operatorname{Fun}(A, B)$, and the map $\beta f: A \rightarrow B^{\mathbb{I}}$ represents an isomorphism in $\operatorname{Fun}(A, B)$ from $f g f$ to $f$. Since $f_{*}$ is an equivalence, we conclude from Remark 1.1.24 that $\mathrm{id}_{A}$ and $g f$ are isomorphic in\\
the quasi-category $\operatorname{Fun}(A, A)$ : explicitly, such an isomorphism may be defined by applying the inverse equivalence $\tilde{h}: \operatorname{Fun}(A, B) \rightarrow \operatorname{Fun}(A, A)$ and composing with the components at $\operatorname{id}_{A}, g f \in \operatorname{Fun}(A, A)$ of the isomorphism $\widetilde{\alpha}: \operatorname{Fun}(A, A) \rightarrow \operatorname{Fun}(A, A)^{\mathbb{I}}$ from $\operatorname{id}_{\operatorname{Fun}(A, A)}$ to $\widetilde{h} f_{*}$. Now by Corollary 1.1.16 this isomorphism is represented by a map $\mathbb{I} \rightarrow \operatorname{Fun}(A, A)$ from $\operatorname{id}_{A}$ to $g f$, which internalizes to a map $\alpha: A \rightarrow A^{\mathbb{I}}$ in $\mathcal{K}$ displayed on the left of (1.2.16).

The converse is easy: the simplicial cotensor construction commutes with $\operatorname{Fun}(X,-)$, so a homotopy equivalence (1.2.16) induces a homotopy equivalence of quasi-categories as in Definition 1.1.23. \(\square\)\\
1.2.17. Lemma. The equivalences in an $\infty$-cosmos are closed under retracts and satisfy the 2 -of-3 property: given a composable pair of functors and their composite, if any two of these are equivalences so is the third.

By the representable definition of equivalences and functoriality, Lemma 1.2.17 follows easily from the corresponding results for equivalences between quasi-categories (see Exercise 1.1.viii). But for sake of completeness, we prove the general cosmological result without relying on this base case, subsuming Exercise 1.1.viii.

Proof. Let $f: A \xrightarrow{口} B$ be an equivalence equipped with the data of (1.2.16) and consider a retract diagram\\
\includegraphics[max width=\textwidth, alt={}, center]{b375083c-8174-49a1-a017-e977744038d8-19_267_389_1046_867}\\
By Lemma 1.2.15, to prove that $h: C \rightarrow D$ is an equivalence, it suffices to construct the data of an inverse homotopy equivalence. To that end define $k: D \rightarrow C$ to be the composite $v g s$ and then observe from the commutative diagrams\\
\includegraphics[max width=\textwidth, alt={}, center]{b375083c-8174-49a1-a017-e977744038d8-19_526_600_1496_427}\\
\includegraphics[max width=\textwidth, alt={}, center]{b375083c-8174-49a1-a017-e977744038d8-19_518_584_1498_1107}\\
that $v^{\mathbb{I}} \alpha u: C \rightarrow C^{\mathbb{I}}$ and $t^{\mathbb{I}} \beta s: D \rightarrow D^{\mathbb{I}}$ define the required homotopy coherent isomorphisms.\\
Via Lemma 1.2.15, the 2-of-3 property for equivalences follows from the fact that the set of isomorphisms in a quasi-category is closed under composition. Homotopy coherent isomorphisms in a quasi-category represent isomorphisms in the homotopy category, whose composite in the homotopy category is then an isomorphism, which can be lifted to a representing homotopy coherent isomorphism by Corollary 1.1.16. ${ }^{17}$ We now apply this to the homotopy coherent isomorphisms in the functor spaces of an $\infty$-cosmos that form part of the data of an equivalence of $\infty$-categories.

\footnotetext{${ }^{17}$ In fact, by Example D.5.5, homotopy coherent isomorphisms can be composed directly, but we do not need this here.
}To prove that equivalences are closed under composition, consider a composable pair of equivalences with their inverse equivalences

$$
A \underset{\stackrel{\sim}{k}}{\stackrel{f}{\rightleftarrows}} B \underset{\stackrel{\sim}{h}}{\stackrel{g}{\underset{\sim}{\sim}}} C
$$

The equivalence data of Lemma 1.2.15 defines isomorphisms $\alpha: \operatorname{id}_{A} \cong k f \in \operatorname{Fun}(A, A)$ and $\gamma: \operatorname{id}_{B} \cong$ $h g \in \operatorname{Fun}(B, B)$, the latter of which whiskers to define $k \gamma f: k f \cong k h g f \in \operatorname{Fun}(B, B)$. Composing these, we obtain an isomorphism $\operatorname{id}_{A} \cong \operatorname{khg} f \in \operatorname{Fun}(A, A)$, witnessing that $k h$ defines a left equivalence inverse of $g f$. The other isomorphism is constructed similarly.

To prove that the equivalences are closed under right cancelation, consider a diagram\\
\includegraphics[max width=\textwidth, alt={}, center]{b375083c-8174-49a1-a017-e977744038d8-20_203_358_747_882}\\
with $k$ an inverse equivalence to $f$ and $\ell$ and inverse equivalence to $g f$. We claim that $f \ell$ defines an inverse equivalence to $g$. One of the required isomorphisms $\mathrm{id}_{C} \cong g f \ell$ is given already. The other is obtained by composing three isomorphisms in $\operatorname{Fun}(B, B)$

$$
\operatorname{id}_{B} \xrightarrow[\bar{z}]{\beta^{-1}} f k \xrightarrow[\bar{z}]{f \delta k} f \ell g f k \xrightarrow[\bar{z}]{f \ell g \beta} f \ell g .
$$

The proof of stability of equivalence under left cancelation is dual. \(\square\)

The trivial fibrations admit a similar characterization as split fiber homotopy equivalences. 1.2.18. Lemma (trivial fibrations split). Every trivial fibration admits a section\\
\includegraphics[max width=\textwidth, alt={}, center]{b375083c-8174-49a1-a017-e977744038d8-20_173_203_1443_964}\\
that defines a split fiber homotopy equivalence\\
\includegraphics[max width=\textwidth, alt={}, center]{b375083c-8174-49a1-a017-e977744038d8-20_265_449_1707_832}\\
and conversely any isofibration that defines a split fiber homotopy equivalence is a trivial fibration.\\
Proof. If $p: E \leadsto B$ is a trivial fibration, then by the final stability property of Lemma 1.2.14, so is $p_{*}: \operatorname{Fun}(X, E) \xrightarrow{\Rightarrow} \operatorname{Fun}(X, B)$ for any $\infty$-category $X$. By Definition 1.1.25, we may solve the lifting problem below-left\\
\includegraphics[max width=\textwidth, alt={}, center]{b375083c-8174-49a1-a017-e977744038d8-20_229_483_2221_483}\\
\includegraphics[max width=\textwidth, alt={}, center]{b375083c-8174-49a1-a017-e977744038d8-20_225_552_2218_1088}\\
to find a map $s: B \rightarrow E$ so that $p s=\mathrm{id}_{B}$, and then solve the lifting problem above-right to construct the desired fibered homotopy. The converse is immediate from Lemma 1.2.15. \(\square\)

A classical construction in abstract homotopy theory proves the following:\\
1.2.19. Lemma (Brown factorization lemma). Any functor $f: A \rightarrow B$ in an $\infty$-cosmos may be factored as an equivalence followed by an isofibration, where this equivalence is constructed as a section of a trivial fibration.\\
\includegraphics[max width=\textwidth, alt={}, center]{b375083c-8174-49a1-a017-e977744038d8-21_214_378_597_872}

(1.2.20)\\
Moreover, $f$ is an equivalence if and only if the isofibration $p$ is a trivial fibration.\\
Proof. The displayed factorization is constructed by the pullback of an isofibration formed by the simplicial cotensor of the inclusion $\mathbb{1}+\mathbb{1} \hookrightarrow \mathbb{I}$ into the $\infty$-category $B$.\\
\includegraphics[max width=\textwidth, alt={}, center]{b375083c-8174-49a1-a017-e977744038d8-21_359_570_1018_780}\\
Note the map $q$ is a pullback of the trivial fibration $\mathrm{ev}_{0}: B^{\mathbb{I}} \xrightarrow{\text { ⇒ }} B$ and is hence a trivial fibration. Its section $s$, constructed by applying the universal property of the pullback to the displayed cone with summit $A$, is thus an equivalence by the 2-of-3 property. Again by 2-of-3, it follows that $f$ is an equivalence if and only if $p$ is. \(\square\)\\
1.2.21. Remark (equivalences satisfy the 2-of-6 property). In fact the equivalences in any $\infty$-cosmos satisfy the stronger 2-of-6 property: for any composable triple of functors\\
\includegraphics[max width=\textwidth, alt={}, center]{b375083c-8174-49a1-a017-e977744038d8-21_306_364_1744_878}\\
if $g f$ and $h g$ are equivalences then $f, g, h$, and $h g f$ are too. An argument of Blumberg and Mandell [20, 6.4] reproduced in Proposition C.1.8 uses Lemmas 1.2.17, 1.2.18, and 1.2.19 to prove that the equivalences have the 2-of-6 property (see Corollary C.1.9).

One of the key advantages of the $\infty$-cosmological approach to abstract category theory is that there are a myriad varieties of "fibered" $\infty$-cosmoi that can be built from a given $\infty$-cosmos, which means that any theorem proven in this axiomatic framework specializes and generalizes to those contexts. The most basic of these derived $\boldsymbol{\infty}$-cosmoi is the $\boldsymbol{\infty}$-cosmos of isofibrations over a fixed base, which\\
we introduce now. Other examples of $\infty$-cosmoi are developed in Chapter 6, once we have a deeper understanding of the cosmological limits of axiom 1.2.1(i).\\
1.2.22. Proposition (sliced $\infty$-cosmoi). For any $\infty$-cosmos $\mathcal{K}$ and any $\infty$-category $B \in \mathcal{K}$ there is an $\infty$-cosmos $\mathcal{K}_{/ B}$ of isofibrations over $B$ whose

\begin{itemize}
  \item[(i)] objects are isofibrations $p: E \rightarrow B$ with codomain $B$
  \item[(ii)] functor spaces, say from $p: E \rightarrow B$ to $q: F \rightarrow B$, are defined by pullback\\
\includegraphics[max width=\textwidth, alt={}, center]{b375083c-8174-49a1-a017-e977744038d8-22_195_707_596_707}\\
and abbreviated to $\operatorname{Fun}_{B}(E, F)$ when the specified isofibrations are clear from context
  \item[(iii)] isofibrations are commutative triangles of isofibrations over $B$\\
\includegraphics[max width=\textwidth, alt={}, center]{b375083c-8174-49a1-a017-e977744038d8-22_140_248_934_937}
  \item[(iv)] terminal object is id: $B \rightarrow B$ and products are defined by the pullback along the diagonal\\
\includegraphics[max width=\textwidth, alt={}, center]{b375083c-8174-49a1-a017-e977744038d8-22_216_346_1173_892}
  \item[(v)] pullbacks and limits of towers of isofibrations are created by the forgetful functor $\mathcal{K}_{/ B} \rightarrow \mathcal{K}$
  \item[(vi)] simplicial cotensor of $p: E \rightarrow B$ with $U \in$ sSet is constructed by the pullback\\
\includegraphics[max width=\textwidth, alt={}, center]{b375083c-8174-49a1-a017-e977744038d8-22_199_335_1535_899}
  \item[(vii)] and in which a map over $B$\\
\includegraphics[max width=\textwidth, alt={}, center]{b375083c-8174-49a1-a017-e977744038d8-22_156_248_1808_938}\\
is an equivalence in the $\infty$-cosmos $\mathcal{K}_{/ B}$ if and only if $f$ is an equivalence in $\mathcal{K}$.
\end{itemize}

Proof. The functor spaces are quasi-categories since axiom 1.2.1(ii) asserts that for any isofibration $q: F \rightarrow B$ in $\mathcal{K}$ the map $q_{*}: \operatorname{Fun}(E, F) \rightarrow \operatorname{Fun}(E, B)$ is an isofibration of quasi-categories. Other parts of this axiom imply that each of the limit constructions - such as the products and cotensors constructed in (iv) and (vi) - define isofibrations over $B$. The closure properties of the isofibrations in $\mathcal{K}_{/ B}$ follow from the corresponding ones in $\mathcal{K}$. The most complicated of these is the Leibniz cotensor stability of the isofibrations in $\mathcal{K}_{/ B}$, which follows from the corresponding property in $\mathcal{K}$, since for a monomorphism of simplicial sets $i: X \hookrightarrow Y$ and an isofibration $r$ over $B$ as in (iii) above, the map $i \widehat{\Pi_{B}} r$ is constructed by pulling back $i \widehat{\pitchfork} r$ along $\Delta: B \rightarrow B^{Y}$.

The fact that the above constructions define simplicially enriched limits in a simplicially enriched slice category are standard from enriched category theory. It remains only to verify that the equivalences in the $\infty$-cosmos of isofibrations are created by the forgetful functor $\mathcal{K}_{/ B} \rightarrow \mathcal{K}$. Suppose first that the map $f$ displayed in (vii) defines an equivalence in $\mathcal{K}$. Then for any isofibration $s: A \rightarrow B$ the induced map on functor spaces in $\mathcal{K}_{/ B}$ is defined by the pullback:\\
\includegraphics[max width=\textwidth, alt={}, center]{b375083c-8174-49a1-a017-e977744038d8-23_268_952_517_585}\\
Since $f$ is an equivalence in $\mathcal{K}$, the map $f_{*}: \operatorname{Fun}(A, E) \rightarrow \operatorname{Fun}(A, F)$ is an equivalence, and so it follows that the induced map on fibers over $s$ is an equivalence as well. ${ }^{18}$

For the converse implication, we appeal to Lemma 1.2.15. If $f: E \rightarrow F$ is an equivalence in $\mathcal{K}_{/ B}$ then it admits a homotopy inverse in $\mathcal{K}_{/ B}$. The inverse equivalence $g: F \rightarrow E$ also defines an inverse equivalence in $\mathcal{K}$ and the required simplicial homotopies in $\mathcal{K}$

$$
E \xrightarrow{\alpha} \mathbb{I} \pitchfork_{B} p \longrightarrow E^{\mathbb{I}} \quad F \xrightarrow{\beta} \mathbb{I} \pitchfork_{B} q \rightarrow F^{\mathbb{I}}
$$

are defined by composing with the top horizontal leg of the pullback defining the cotensor in $\mathcal{K}_{/ B}$. \(\square\)

As mentioned in Digression 1.2.13, many of the $\infty$-cosmoi we encounter "in the wild" satisfy an additional axiom. Note, however, that this axiom is not inherited by the sliced $\boldsymbol{\infty}$-cosmoi of Proposition 1.2.22, which is one of the reasons it was not included in Definition 1.2.1.\\
1.2.23. Definition (cartesian closed $\infty$-cosmoi). An $\infty$-cosmos $\mathcal{K}$ is cartesian closed if the product bifunctor $-\times-: \mathcal{K} \times \mathcal{K} \rightarrow \mathcal{K}$ extends to a simplicially enriched two-variable adjunction

$$
\operatorname{Fun}(A \times B, C) \cong \operatorname{Fun}\left(A, C^{B}\right) \cong \operatorname{Fun}\left(B, C^{A}\right)
$$

in which the right adjoints $(-)^{A}: \mathcal{K} \rightarrow \mathcal{K}$ preserve isofibrations for all $A \in \mathcal{K}$.\\
For instance, the $\infty$-cosmos of quasi-categories is cartesian closed, with the exponentials defined as (special cases of) simplicial cotensors. This is one of the reasons that we use the same notation for cotensor and for exponential. ${ }^{19}$ Note in this case the functor spaces and the exponentials coincide. The same is true for the cartesian closed $\infty$-cosmoi of categories and of Kan complexes. In general, the functor space from $A$ to $B$ is the "underlying quasi-category" of the exponential $B^{A}$ whenever it exists (see Remark 1.3.11).\\
1.2.24. Example ( $\infty$-cosmoi of $(\infty, 1)$-categories; §E.2). The following models of $(\infty, 1)$-categories define cartesian closed $\infty$-cosmoi:

\footnotetext{${ }^{18}$ The stability of equivalences between isofibrations under pullback can be proven either as a consequence of Lemmas 1.2.14 and 1.2.19 using standard techniques from simplicial homotopy theory (see Lemma C.1.11) or by arguing 2-categorically (see Proposition 3.3.4).\\
${ }^{19}$ Other advantages of this convenient notational conflation are discussed in §2.3 and in Proposition 10.3.5.
}
\begin{itemize}
  \item[(i)] Rezk's complete Segal spaces define the objects of an $\infty-\cos \operatorname{mos} \mathcal{C S S}$, in which the isofibrations, equivalences, and trivial fibrations are the corresponding classes of the model structure of [100]. ${ }^{20}$
  \item[(ii)] The Segal categories defined by Dwyer, Kan, and Smith [38] and developed by Hirschowitz and Simpson [56] define the objects of an $\infty$-cosmos Segal, in which the isofibrations, equivalences, and trivial fibrations are the corresponding classes of the model structure of [13, 90]. ${ }^{21}$
  \item[(iii)] The 1-complicial sets of [129], equivalently the "naturally marked quasi-categories" of [78], define the objects of an $\infty$-cosmos 1-Comp in which the isofibrations, equivalences, and trivial fibrations are the corresponding classes of the model structure from either of these sources.
\end{itemize}


In §E.3, we show that certain models of $(\infty, n)$-categories or even $(\infty, \infty)$-categories define $\infty$-cosmoi: $n$-quasi-categories, $\Theta_{n}$-spaces, iterated complete Segal spaces, and $n$-complicial sets.\\
1.2.25. Definition (co-dual $\infty$-cosmoi). There is an identity-on-objects involutive functor $(-)^{\circ}: \boldsymbol{\Delta} \rightarrow$ $\boldsymbol{\Delta}$ that reverses the ordering of the elements in each ordinal $[n] \in \boldsymbol{\Delta}$. In the notation of 1.1.1, the functor (-)° sends a face map $\delta^{i}:[n-1] \gtrdot[n]$ to the face map $\delta^{n-i}:[n-1] \succ[n]$ and sends the degeneracy map $\sigma^{i}:[n+1] \rightarrow[n]$ to the degeneracy map $\sigma^{n-i}:[n+1] \rightarrow[n]$. Precomposition with this involutive automorphism induces an involution (-) ${ }^{\mathrm{op}}: s \mathcal{S} e t \rightarrow s \mathcal{S} e t$ that sends a simplicial set $X$ to its opposite simplicial set $X^{\mathrm{op}}$, with the orientation of the vertices in each simplex reversed. This construction preserves all conical limits and colimits and induces an isomorphism $\left(Y^{X}\right)^{\mathrm{op}} \cong\left(Y^{\mathrm{op}}\right)^{X^{\mathrm{op}}}$ on exponentials.

For any $\infty$-cosmos $\mathcal{K}$, there is a dual $\infty$-cosmos $\mathcal{K}^{\mathrm{co}}$ with the same objects but with functor spaces defined by:

$$
\operatorname{Fun}_{\mathcal{K}^{\mathrm{co}}}(A, B):=\operatorname{Fun}_{\mathcal{K}}(A, B)^{\mathrm{op}} .
$$

The isofibrations, equivalences, and trivial fibrations in $\mathcal{K}^{\mathrm{co}}$ coincide with those of $\mathcal{K}$.\\
Conical limits in $\mathcal{K}^{\mathrm{c}^{\mathrm{o}}}$ coincide with those in $\mathcal{K}$, while the cotensor of $A \in \mathcal{K}$ with $U \in s \mathcal{S} e t$ is defined to be $A^{U^{\mathrm{op}}}$.

A 2-categorical justification for this notation is given in Exercise 1.4.ii.\\
1.2.26. Definition (discrete $\infty$-categories). An $\infty$-category $E$ in an $\infty$-cosmos $\mathcal{K}$ is discrete just when for all $X \in \mathcal{K}$ the functor space $\operatorname{Fun}(X, E)$ is a Kan complex.

In the $\infty$-cosmos of quasi-categories, the discrete $\infty$-categories are exactly the Kan complexes. Similarly, in the $\infty$-cosmoi of Example 1.2.24 whose $\infty$-categories are $(\infty, 1)$-categories in some model, the discrete $\boldsymbol{\infty}$-categories are the $\boldsymbol{\infty}$-groupoids. Importantly for what follows, the discrete $\boldsymbol{\infty}$-categories can be characterized "internally" to the $\infty$-cosmos as follows:\\
1.2.27. Lemma. An $\infty$-category $E$ is discrete if and only if $E^{\mathbb{I}} \leadsto E^{2}$ is a trivial fibration.

\footnotetext{${ }^{20}$ Warning: the model category of complete Segal spaces is enriched over simplicial sets in two distinct "directions" - one enrichment makes the simplicial set of maps between two complete Segal spaces into a Kan complex that probes the "spacial" structure while another enrichment makes the simplicial set of maps into a quasi-category that probes the "categorical" structure [64]. It is this latter enrichment that we want.\\
${ }^{21}$ Here we reserve the term "Segal category" for those simplicial objects with a discrete set of objects that are Reedy fibrant and satisfy the Segal condition. The traditional definition does not include the Reedy fibrancy condition because it is not satisfied by the simplicial object defined as the nerve of a Kan complex enriched category. Since Kan complex enriched categories are not among our preferred models of $(\infty, 1)$-categories this does not bother us.
}Proof. By Definition 1.2.2, the isofibration $E^{\mathbb{I}} \rightarrow E^{2}$ is a trivial fibration if and only if for all $\infty$-categories $X$ the induced map on functor spaces

$$
\begin{aligned}
\operatorname{Fun}\left(X, E^{\mathbb{I}}\right) & \longrightarrow \operatorname{Fun}\left(X, E^{2}\right) \\
\text { RII } & \text { RII } \\
\operatorname{Fun}(X, E)^{\mathbb{I}} & \longrightarrow \operatorname{Fun}(X, E)^{2}
\end{aligned}
$$

is a trivial fibration of quasi-categories. Via the universal property of the simplicial cotensor, Lemma 1.1.30 tells us that this map is a trivial fibration if and only if $\operatorname{Fun}(X, E)$ is a Kan complex. \(\square\)

The reader may check that the discrete $\infty$-categories in any $\infty$-cosmos assemble into an $\infty$-cosmos $\mathcal{K}^{\sim}$. A proof appears in Proposition 6.1.6 where general techniques for producing new $\infty$-cosmoi from given ones are developed.

\section*{Exercises.}
1.2.i. Exercise. Define an equivalence between the categories of:

\begin{itemize}
  \item[(i)] simplicial categories, as in (1.2.5), and
  \item[(ii)] categories enriched over simplicial sets.
\end{itemize}

1.2.ii. Exercise. Elaborate on the proof of Proposition 1.2.10 by proving that the simplicially enriched category QCat admits conical products satisfying the universal property of Digression 1.2.6. That is:

\begin{itemize}
  \item[(i)] Define the cartesian product $A \times B$ and the projection maps $\pi_{A}: A \times B \rightarrow A$ and $\pi_{B}: A \times B \rightarrow$ $B$ for a pair of quasi-categories $A$ and $B$ and prove that this data satisfies the usual (unenriched) universal property.
  \item[(ii)] Given another quasi-category $X$, use (i) and the Yoneda lemma to show that the projection maps induce an isomorphism of quasi-categories
$$
(A \times B)^{X} \xrightarrow{\simeq} A^{X} \times B^{X} .
$$
  \item[(iii)] Explain how this relates to the universal property of Digression 1.2.6.
  \item[(iv)] Express the usual 1-categorical universal property of (i) as the "0-dimensional aspect" of the universal property of (ii).
\end{itemize}

1.2.iii. Exercise. Prove that any object in an $\infty$-cosmos has a path object\\
\includegraphics[max width=\textwidth, alt={}, center]{b375083c-8174-49a1-a017-e977744038d8-25_198_437_1791_843}\\
constructed by cotensoring with the free-living isomorphism.\\
1.2.iv. Exercise. Show that if $\mathcal{K}$ is a cartesian closed $\infty$-cosmos then $\mathcal{K}^{\mathrm{co}}$ is as well.\\
1.2.v. Exercise (6.1.6). Use Proposition 1.2.12 to show that the discrete $\infty$-categories in any $\infty$-cosmos define an $\infty$-cosmos whose functor spaces are all Kan complexes.

\subsection*{1.3. Cosmological Functors}
Certain "right adjoint type" constructions define maps between $\infty$-cosmoi that preserve all of the structures axiomatized in Definition 1.2.1. The simple observation that such constructions define cosmological functors between $\infty$-cosmoi streamlines many proofs.\\
1.3.1. Definition (cosmological functor). A cosmological functor is a simplicial functor (see Definition A.2.6) between $\infty$-cosmoi that preserves the specified isofibrations and all of the cosmological limits.

In general, cosmological functors preserve any $\infty$-categorical notion that can be characterized internally to the $\infty$-cosmos - for instance, as a map equipped with additional structure - as opposed to externally - for instance, by a statement that involves a universal or existential quantifier. For example, the equivalences in an $\infty$-cosmos are characterized externally in Definition 1.2.2, which might lead one to suspect that a nonsurjective cosmological functor could fail to preserve them. However, Lemma 1.2.15 characterizes equivalences in terms of the presence of structures defined internally to an $\infty$-cosmos, so as a result:\\
1.3.2. Lemma. Any cosmological functor also preserves equivalences and trivial fibrations.

Proof. By Lemma 1.2.15 the equivalences in an $\infty$-cosmos coincide with the "homotopy equivalences" defined by cotensoring with the free-living isomorphism. Since a cosmological functor preserves simplicial cotensors, it preserves the data displayed in (1.2.16) and hence carries equivalences to equivalences. The preservation of trivial fibrations follows. \(\square\)\\
1.3.3. Remark. Similarly, arguing from Definition 1.2.26 it would not be clear whether cosmological functors preserve discrete $\infty$-categories, but using the internal characterization of Lemma 1.2.27 - an $\infty$-category $A$ is discrete if and only if $A^{\mathbb{I}} \xrightarrow{ } A^{2}$ is a trivial fibration - this follows from the fact that cosmological functors preserve simplicial cotensors and trivial fibrations.

We now demonstrate that cosmological functors are abundant:\\
1.3.4. Proposition. The following constructions define cosmological functors for any $\infty$-cosmos $\mathcal{K}$ :

\begin{itemize}
  \item[(i)] The functor space $\operatorname{Fun}(X,-): \mathcal{K} \rightarrow Q$ Cat, for any $\infty$-category $X$.
  \item[(ii)] The underlying quasi-category functor
$$
(-)_{0}:=\operatorname{Fun}(1,-): \mathcal{K} \rightarrow Q C a t,
$$
specializing (i) to the terminal $\infty$-category 1.
  \item[(iii)] The simplicial cotensor $(-)^{U}: \mathcal{K} \rightarrow \mathcal{K}$, for any simplicial set $U$.
  \item[(iv)] The exponential $(-)^{A}: \mathcal{K} \rightarrow \mathcal{K}$, for any $\infty$-category $A$ in a cartesian closed $\infty$-cosmos $\mathcal{K}$.
  \item[(v)] Pullback of isofibrations $f^{*}: \mathcal{K}_{/ B} \rightarrow \mathcal{K}_{/ A}$ along any functor $f: A \rightarrow B$ in an $\infty$-cosmos $\mathcal{K}$.
  \item[(vi)] Moreover, for any cosmological functor $F: \mathcal{K} \rightarrow \mathcal{L}$ and any $\infty$-category $A \in \mathcal{K}$, the induced map on slices $F: \mathcal{K}_{/ A} \rightarrow \mathcal{L}_{/ F A}$ defines a cosmological functor.
\end{itemize}

Proof. The first four of these statements are nearly immediate, the preservation of isofibrations being asserted explicitly as a hypothesis in each case and the preservation of limits following from familiar arguments.

For (v), pullback in an $\infty$-cosmos $\mathcal{K}$ is a simplicially enriched limit construction; one consequence of this is that $f^{*}: \mathcal{K}_{/ B} \rightarrow \mathcal{K}_{/ A}$ defines a simplicial functor. The action of the functor $f^{*}$ on a 0 -arrow\\
$g$ in $\mathcal{K}_{/ B}$ is also defined by a pullback square: since the front and back squares in the displayed diagram are pullbacks the top square is as well\\
\includegraphics[max width=\textwidth, alt={}, center]{b375083c-8174-49a1-a017-e977744038d8-27_284_452_371_831}\\
Since isofibrations are stable under pullback, it follows that $f^{*}: \mathcal{K}_{/ B} \rightarrow \mathcal{K}_{/ A}$ preserves isofibrations. It remains to prove that this functor preserves the simplicial limits constructed in Proposition 1.2.22, which is fundamentally a consequence of the commutativity of limit constructions. In each case, this can be verified explicitly. We illustrate this computation for simplicial cotensors by constructing the commutative cube:\\
\includegraphics[max width=\textwidth, alt={}, center]{b375083c-8174-49a1-a017-e977744038d8-27_392_678_914_727}\\
Since the front, back, and right faces are pullbacks, the left is as well.\\
The final statement (vi) is left as Exercise 1.3.i. \(\square\)\\
1.3.5. Example. By Propositions 1.2.11 and 1.2.12, the full subcategory inclusions Cat $\hookrightarrow$ QCat and Kan ↪ QCat both define cosmological functors (see also Lemma 6.1.4). These cosmological embeddings explicate the intuition that the formal category theory of 1-categories or of $\infty$-groupoids can be recovered as a special case of the formal category theory of $(\infty, 1)$-categories.

Non-examples of cosmological functors are also instructive:\\
1.3.6. Non-Example. The forgetful functor $\mathcal{K}_{/ B} \rightarrow \mathcal{K}$ is simplicial and preserves isofibrations but does not define a cosmological functor, failing to preserve cotensors and products. However, by Proposition 1.3.4(v), its right adjoint $-\times B: \mathcal{K} \rightarrow \mathcal{K}_{/ B}$ does define a cosmological functor.\\
1.3.7. Non-Example. The cosmological embedding Kan →QCat has a right adjoint $(-)^{\sim}:$ QCat → Kan that carries each quasi-category to its " $\infty$-groupoid core" or maximal sub Kan complex, the simplicial subset containing those $n$-simplices whose edges are all isomorphisms. This core functor preserves isofibrations and 1-categorical limits but is not cosmological since it is not simplicially enriched: any functor $K \rightarrow Q$ whose domain is a Kan complex and whose codomain is a quasi-category factors through the inclusion $Q^{\simeq} \hookrightarrow Q$ via a unique map $K \rightarrow Q^{\simeq}$ but in general $\operatorname{Fun}(K, Q) \neq$ $\operatorname{Fun}\left(K, Q^{\simeq}\right)$, since a natural transformation $K \times \Delta[1] \rightarrow Q$ only factors through $Q^{\simeq} \hookrightarrow Q$ in the case where its components are invertible (see Lemma 12.1.12 however).

Certain cosmological functors are especially well-behaved:\\
1.3.8. Definition (cosmological biequivalence). A cosmological functor defines a cosmological biequivalence $F: \mathcal{K} \leadsto \mathcal{L}$ if it additionally

\begin{itemize}
  \item[(i)] is essentially surjective on objects up to equivalence: for all $C \in \mathcal{L}$ there exists $A \in \mathcal{K}$ so that $F A \simeq C$ and
  \item[(ii)] defines a local equivalence: for all $A, B \in \mathcal{K}$, the action of $F$ on functor spaces defines an equivalence of quasi-categories
$$
\operatorname{Fun}(A, B) \longrightarrow \operatorname{Fun}(F A, F B) .
$$
\end{itemize}

Cosmological biequivalences are studied more systematically in Chapter 10, where we think of them as "change-of-model" functors. Crucially for our proof of the "model independence" of $(\infty, 1)$-category theory in Chapter 11, there are a variety of cosmological biequivalences between the $\infty$-cosmoi of $(\infty, 1)$-categories:\\
1.3.9. Example (biequivalences between $\infty$-cosmoi of $(\infty, 1)$-categories; §E.2).

\begin{itemize}
  \item[(i)] The underlying quasi-category functors defined on the $\infty$-cosmoi of complete Segal spaces, Segal categories, and 1-complicial sets
$$
\text { CSS } \xrightarrow[\sim]{(-)_{0}} \text { QCat } \quad \text { Segal } \xrightarrow[\sim]{(-)_{0}} \text { QCat } \quad 1-\text { Comp } \xrightarrow[\sim]{(-)_{0}} \text { QCat }
$$
are all biequivalences. In the first two cases these are defined by "evaluating at the 0th row" and in the last case this is defined by "forgetting the markings."
  \item[(ii)] There are also cosmological biequivalences nerve: $Q$ Cat $\xrightarrow{\sim} C S \mathcal{S}$ and nerve: $Q$ Cat $\xrightarrow{\sim}$ Segal defined by Joyal and Tierney [64].
  \item[(iii)] The functor disc: $C S S \bumpeq$ Segal defined by Bergner [14] that "discretizes" a complete Segal spaces also defines a cosmological biequivalence.
  \item[(iv)] Another cosmological biequivalence $(-)^{\natural}: Q$ Cat $\xrightarrow{\sim} 1-$ Comp that gives each quasi-category its "natural marking."
\end{itemize}

In terminology justified by Proposition 10.2.1:\\
1.3.10. Definition. An $\infty$-cosmos $\mathcal{K}$ is an $\infty$-cosmos of $(\infty, 1)$-categories just when the underlying quasi-category $(-)_{0}: \mathcal{K} \rightarrow$ QCat is a cosmological biequivalence.\\
1.3.11. Remark. The underlying quasi-category functor $(-)_{0}: \mathcal{K} \rightarrow Q C$ at carries the internal homs of a cartesian closed $\infty$-cosmos $\mathcal{K}$ to the corresponding functor spaces: for any $\infty$-categories $A$ and $B$ in $\mathcal{K}$, we have

$$
\left(B^{A}\right)_{0}:=\operatorname{Fun}\left(1, B^{A}\right) \cong \operatorname{Fun}(A, B) .
$$

In the case where the $\infty$-cosmos $\mathcal{K}$ is biequivalent to $Q$ Cat, we see in Chapters 10 and 11 that this entails no essential loss of categorical information.

Cosmological biequivalences not only preserve equivalences but also reflect and create them.\\
1.3.12. Lemma. Let $F: \mathcal{K} \xrightarrow{\boldsymbol{L}}$ be a cosmological biequivalence. Then:

\begin{itemize}
  \item[(i)] A functor $f: A \rightarrow B$ between $\infty$-categories in $\mathcal{K}$ is an equivalence if and only if $F f: F A \rightarrow F B$ is an equivalence in $\mathcal{L}$.
  \item[(ii)] A pair of $\infty$-categories in $\mathcal{K}$ are equivalent if and only if their images in $\mathcal{L}$ are equivalent.
\end{itemize}

Proof. Lemma 1.3.2 implies that cosmological functors preserve equivalences and thus also the existence of an equivalence between a pair of $\infty$-categories in $\mathcal{K}$. To see that equivalences are also\\
reflected, suppose $f: A \rightarrow B$ is a functor in $\mathcal{K}$ with the property that $F f: F A \xrightarrow{\Rightarrow} F B$ is an equivalence in $\mathcal{L}$. Now for any $\infty$-category $X$, simplicial functoriality provides a commutative diagram\\
\includegraphics[max width=\textwidth, alt={}, center]{b375083c-8174-49a1-a017-e977744038d8-29_215_569_367_778}\\
so from the 2-of-3 property we conclude that $f_{*}: \operatorname{Fun}(X, A) \xrightarrow{\sim} \operatorname{Fun}(X, B)$ is an equivalence, proving that $f$ is an equivalence in $\mathcal{K}$.

To see that equivalences are created, suppose now that $A$ and $B$ are $\infty$-categories in $\mathcal{K}$ equipped with an equivalence:\\
\includegraphics[max width=\textwidth, alt={}, center]{b375083c-8174-49a1-a017-e977744038d8-29_317_1231_820_449}\\
in $\mathcal{L}$. Since $\operatorname{Fun}(A, B) \xrightarrow{ } \operatorname{Fun}(F A, F B)$ and $\operatorname{Fun}(B, A) \xrightarrow{\sim} \operatorname{Fun}(F B, F A)$ are equivalences of quasi-categories the induced functors between their homotopy categories $\mathrm{h}(\operatorname{Fun}(A, B)) \leadsto \mathrm{h}(\operatorname{Fun}(F A, F B))$ and $\mathrm{h}(\operatorname{Fun}(B, A)) \xrightarrow{\sim} \mathrm{h}(\operatorname{Fun}(F B, F A))$ are equivalences of categories, by Remark 1.1.24, and in particular essentially surjective. So we may lift $\tilde{f}$ and $\tilde{g}$ to functors $f: A \rightarrow B$ and $g: B \rightarrow A$ inh (Fun( $A, B)$ ) and $\mathrm{h}(\operatorname{Fun}(B, A))$, respectively, so that $F f \cong \widetilde{f}$ and $F g \cong \tilde{g}$. The commutative diagram of quasi-categories\\
\includegraphics[max width=\textwidth, alt={}, center]{b375083c-8174-49a1-a017-e977744038d8-29_194_828_1433_648}\\
induces a commutative diagram between their homotopy categories. In particular, by applying the composition bifunctor to the isomorphisms $F f \cong \tilde{f}$ and $F g \cong \tilde{g}$, we see that

$$
F\left(\operatorname{id}_{A}\right)=\operatorname{id}_{F A} \cong \tilde{g} \circ \tilde{f} \cong F g \circ F f=F(g \circ f)
$$

in h( $\operatorname{Fun}(F A, F A)$ ). By fully faithfulness of $\mathrm{h}(\operatorname{Fun}(A, A)) \xrightarrow{\sim} \mathrm{h}(\operatorname{Fun}(F A, F A))$, this isomorphism lifts to an isomorphism $\mathrm{id}_{A} \cong g \circ f$ in $\mathrm{h}(\operatorname{Fun}(A, A))$. By Corollary 1.1.16, this isomorphism can be represented by a homotopy coherent isomorphism $\mathbb{I} \rightarrow \operatorname{Fun}(A, A)$, which internalizes to define a map $\alpha: A \rightarrow A^{\mathbb{I}}$ as required. The construction of the homotopy coherent isomorphism $\beta: B \rightarrow B^{\mathbb{I}}$ from $f \circ g$ to $\operatorname{id}_{B}$ proceeds similarly. \(\square\)

The proof of the creation of equivalences in Lemma 1.3.12 is surprisingly delicate, passing to the homotopy categories of the functor spaces to avoid lifting and composing homotopy coherent isomorphisms; an argument along those lines is also possible, and left to the reader as Exercise 1.3.ii. The next section provides context for the argument just given by introducing the homotopy 2-category of an $\infty$-cosmos. The reader is then invited to revisit the creation of equivalences in Exercise \href{http://1.4.vi}{1.4.vi}.

\section*{Exercises.}
1.3.i. Exercise. Prove that for any cosmological functor $F: \mathcal{K} \rightarrow \mathcal{L}$ and any $A \in \mathcal{K}$, the induced map $F: \mathcal{K}_{/ A} \rightarrow \mathcal{L}_{/ F A}$ defines a cosmological functor.\\
1.3.ii. Exercise. Sketch a proof that cosmological biequivalences create equivalences between $\infty$-categor ies without passing to homotopy categories, by lifting and composing the homotopy coherent isomorphisms given as part of the data of the hypothesized equivalences.\\
1.3.iii. Exercise. Suppose $F: \mathcal{K} \rightarrow \mathcal{L}, G: \mathcal{L} \rightarrow \mathcal{M}$, and $H: \mathcal{M} \rightarrow \mathcal{N}$ are cosmological functors, and assume that $G F$ and $H G$ are cosmological biequivalences. Show that $F, G, H$, and $H G F$ are cosmological biequivalences.

\subsection*{1.4. The Homotopy 2-Category}
Small 1-categories define the objects of a strict 2-category ${ }^{22} C a t$ of categories, functors, and natural transformations. Many basic categorical notions - those defined in terms of categories, functors, and natural transformations - can be defined internally to the 2-category Cat. This suggests a natural avenue for generalization: reinterpreting these same definitions in a generic 2-category using its objects in place of small categories, its 1-cells in place of functors, and its 2-cells in place of natural transformations.

In Chapter 2, we develop a significant portion of the theory of $\infty$-categories in any fixed $\infty$ cosmos following exactly this outline, working internally to a 2-category that we refer to as the homotopy 2-category that we associate to any $\infty$-cosmos. The homotopy 2-category of an $\infty$-cosmos is a quotient of the full $\infty$-cosmos, replacing each quasi-categorical functor space by its homotopy category. Surprisingly, this rather destructive quotienting operation preserves quite a lot of information. Indeed, essentially all of the development of the theory of $\infty$-categories in Part I takes place in the homotopy 2 -category of an $\infty$-cosmos. This said, we caution the reader against becoming overly seduced by homotopy 2-categories, which are more of a technical convenience for reducing the complexity of our arguments than a fundamental notion of $\infty$-category theory.

The homotopy 2-category for the $\infty$-cosmos of quasi-categories was first introduced by Joyal in his work on the foundations of quasi-category theory [63].\\
1.4.1. Definition (homotopy 2-category). Let $\mathcal{K}$ be an $\infty$-cosmos. Its homotopy 2-category is the 2-category $\mathfrak{h} \mathcal{K}$ whose

\begin{itemize}
  \item[-] objects are the objects $A, B$ of $\mathcal{K}$, i.e., the $\infty$-categories;
  \item[-] 1-cells $f: A \rightarrow B$ are the 0-arrows in the functor space $\operatorname{Fun}(A, B)$, i.e., the $\infty$-functors; and
  \item[-] 2-cells $A \xrightarrow[g]{\underset{g}{f}} B$ are homotopy classes of 1-simplices in $\operatorname{Fun}(A, B)$, which we call $\infty$-natural transformations.
\end{itemize}

\footnotetext{${ }^{22}$ Appendix B introduces 2-categories and 2-functors, reviewing the 2-category theory needed here. Succinctly, in parallel with Digression 1.2.4, 2-categories (see Definition B.1.1) can be understood equally as:

\begin{itemize}
  \item[-] "two-dimensional" categories, with objects; 1-cells, whose boundary are given by a pair of objects; and 2-cells, whose boundary are given by a parallel pair of 1-cells between a pair of objects - together with partially defined composition operations governed by this boundary data
  \item[-] or as categories enriched over Cat.
\end{itemize}
}Put another way $\mathfrak{h} \mathcal{K}$ is the 2-category with the same objects as $\mathcal{K}$ and with hom-categories defined by

$$
\mathrm{hFun}(A, B):=\mathrm{h}(\operatorname{Fun}(A, B)),
$$

that is, $\mathrm{hFun}(A, B)$ is the homotopy category of the quasi-category $\operatorname{Fun}(A, B)$.\\
The underlying category of a 2-category is defined by simply forgetting its 2-cells. Note that an $\infty$-cosmos $\mathcal{K}$ and its homotopy 2-category $\mathfrak{h} \mathcal{K}$ share the same underlying category $\mathcal{K}_{0}$ of $\infty$-categories and $\infty$-functors in $\mathcal{K}$.\\
1.4.2. Digression (change of base, §A.7). The homotopy category functor preserves finite products, as of course does its right adjoint. It follows that the adjunction of Proposition 1.1.11 induces a change-of-base adjunction\\
\includegraphics[max width=\textwidth, alt={}, center]{b375083c-8174-49a1-a017-e977744038d8-31_128_490_777_816}\\
whose left and right adjoints change the enrichment by applying the homotopy category functor or the nerve functor to the hom objects of the enriched category. Here 2-Cat and sSet-Cat can each be understood as 2-categories - of enriched categories, enriched functors, and enriched natural transformations - and both change of base constructions define 2-functors (see Propositions A.7.3 and A.7.5). Since the nerve embedding is fully faithful, 2-categories can be identified as a full subcategory comprised of those simplicial categories whose hom spaces are nerves of categories.

The proof of Lemma 1.3.12 uses an observation worth highlighting:\\
1.4.3. Lemma.

\begin{itemize}
  \item[(i)] Every 2-cell $A \underset{g}{\stackrel{\Downarrow}{g}} B$ in the homotopy 2-category of an $\infty$-cosmos is represented by a map of quasi-categories as below-left or equivalently by a functor as below-right\\
\includegraphics[max width=\textwidth, alt={}, center]{b375083c-8174-49a1-a017-e977744038d8-31_162_912_1556_605}\\
and two such maps represent the same 2 -cell if and only if they are homotopic as 1 -simplices in $\operatorname{Fun}(A, B)$.
  \item[(ii)] Every invertible 2-cell $A \xrightarrow[g]{\cong} B$ in the homotopy 2-category of an $\infty$-cosmos is represented by a map of quasi-categories as below-left or equivalently by a functor as below-right\\
\includegraphics[max width=\textwidth, alt={}, center]{b375083c-8174-49a1-a017-e977744038d8-31_162_901_1995_612}\\
and two such maps represent the same invertible 2-cell if and only if their common restrictions along $2 \hookrightarrow \mathbb{I}$ are homotopic as 1 -simplices in $\operatorname{Fun}(A, B)$.
\end{itemize}

The notion of homotopic 1-simplices referenced here is defined in Lemma 1.1.9. Since the 2-cells in the homotopy 2-category are referred to as $\infty$-natural transformations, we refer to the invertible 2-cells in the homotopy 2-category as $\infty$-natural isomorphisms.

Proof. The statement (i) records the definition of the 2-cells in the homotopy 2-category and the universal property (1.2.7) of the simplicial cotensor. For (ii), a 2-cell in the homotopy 2-category is invertible if and only if it defines an isomorphism in the appropriate hom-category $\mathrm{hFun}(A, B)$. By Corollary 1.1.16 it follows that each invertible 2-cell $\alpha$ is represented by a homotopy coherent isomorphism $\alpha: \mathbb{I} \rightarrow \operatorname{Fun}(A, B)$, which similarly internalizes to define a functor ${ }^{「} \alpha^{\urcorner}: A \rightarrow B^{\mathbb{I}}$. \(\square\)

An upshot of Digression 1.4.2 is that change of base is an operation that applies to enriched functors as well as enriched categories, as can be directly verified in the case of greatest interest.\\
1.4.4. Lemma. Any simplicial functor $F: \mathcal{K} \rightarrow \mathcal{L}$ between $\infty$-cosmoi induces a 2 -functor $F: \mathfrak{h} \mathcal{K} \rightarrow \mathfrak{h} \mathcal{L}$ between their homotopy 2-categories.

Proof. The action of the induced 2-functor $F: \mathfrak{h} \mathcal{K} \rightarrow \mathfrak{h} \mathcal{L}$ on objects and 1-cells is given by the corresponding action of $F: \mathcal{K} \rightarrow \mathcal{L}$; recall an $\infty$-cosmos and its homotopy 2-category have the same underlying 1-category. Each 2-cell in $\mathfrak{h K}$ is represented by a 1-simplex in $\operatorname{Fun}(A, B)$ which is mapped via\\
\includegraphics[max width=\textwidth, alt={}, center]{b375083c-8174-49a1-a017-e977744038d8-32_203_716_950_704}\\
to a 1-simplex representing a 2-cell in $\mathfrak{h} \mathcal{L}$. Since the action $F: \operatorname{Fun}(A, B) \rightarrow \operatorname{Fun}(F A, F B)$ on functor spaces defines a morphism of simplicial sets, it preserves faces and degeneracies. In particular, homotopic 1-simplices in $\operatorname{Fun}(A, B)$ are carried to homotopic 1-simplices in $\operatorname{Fun}(F A, F B)$ so the action on 2-cells just described is well-defined. The 2-functoriality of these mappings follows from the simplicial functoriality of the original mapping. \(\square\)

We now begin to relate the simplicially enriched structures of an $\infty$-cosmos to the 2-categorical structures in its homotopy 2-category by proving that homotopy 2-categories inherit products from their $\infty$-cosmoi that satisfy a 2-categorical universal property. To illustrate, recall that the terminal $\infty$-category $1 \in \mathcal{K}$ has the universal property $\operatorname{Fun}(X, 1) \cong \mathbb{1}$ for all $X \in \mathcal{K}$. Applying the homotopy category functor we see that $1 \in \mathfrak{h} \mathcal{K}$ has the universal property $\mathrm{hFun}(X, 1) \cong \mathbb{1}$ for all $X \in \mathfrak{h} \mathcal{K}$, which is expressed by saying that the $\infty$-category 1 defines a 2 -terminal object in the homotopy 2 -category. This 2-categorical universal property has both a 1-dimensional and a 2-dimensional aspect. Since $\mathrm{hFun}(X, 1) \cong \mathbb{1}$ is a category with a single object, there exists a unique morphism $X \rightarrow 1$ in $\mathcal{K}$, and since $\mathrm{hFun}(X, 1) \cong \mathbb{1}$ has only a single morphism, the only 2-cells in $\mathfrak{h} \mathcal{K}$ with codomain 1 are identities.

\subsection*{1.4.5. Proposition (cartesian (closure)).}
\begin{itemize}
  \item[(i)] The homotopy 2-category of any $\infty$-cosmos has 2-categorical products.
  \item[(ii)] The homotopy 2-category of a cartesian closed $\infty$-cosmos is cartesian closed as a 2-category.
\end{itemize}

Proof. While the functor h: sSet → Cat only preserves finite products, the restricted functor h: QCat → Cat preserves all products on account of the simplified description of the homotopy category of a quasi-category given in Lemma 1.1.12. Thus for any set $I$ and family of $\infty$-categories $\left(A_{i}\right)_{i \in I}$ in $\mathcal{K}$, the homotopy category functor carries the isomorphism of functor spaces to an isomorphism of hom-categories

$$
\operatorname{Fun}\left(X, \prod_{i \in I} A_{i}\right) \xrightarrow{\cong} \prod_{i \in I} \operatorname{Fun}\left(X, A_{i}\right) \quad \stackrel{\mathrm{h}}{\longrightarrow} \quad \mathrm{hFun}\left(X, \prod_{i \in I} A_{i}\right) \xrightarrow{\cong} \prod_{i \in I} \mathrm{hFun}\left(X, A_{i}\right) .
$$

This proves that the homotopy 2-category $\mathfrak{h} \mathcal{K}$ has products whose universal properties have both a 1-and 2-dimensional component, as described in the empty case for terminal objects above.

If $\mathcal{K}$ is a cartesian closed $\infty$-cosmos, then for any triple of $\infty$-categories $A, B, C \in \mathcal{K}$ there exist exponential objects $C^{A}, C^{B} \in \mathcal{K}$ characterized by natural isomorphisms

$$
\operatorname{Fun}(A \times B, C) \cong \operatorname{Fun}\left(A, C^{B}\right) \cong \operatorname{Fun}\left(B, C^{A}\right) .
$$

Passing to homotopy categories we have natural isomorphisms

$$
\mathrm{hFun}(A \times B, C) \cong \mathrm{hFun}\left(A, C^{B}\right) \cong \mathrm{hFun}\left(B, C^{A}\right),
$$

which demonstrates that $\mathfrak{h} \mathcal{K}$ is cartesian closed as a 2-category: functors $A \times B \rightarrow C$ transpose to define functors $A \rightarrow C^{B}$ and $B \rightarrow C^{A}$, and natural transformations transpose similarly. \(\square\)

There is a standard definition of isomorphism between two objects in any 1-category, preserved by any functor. Similarly, there is a standard definition of equivalence between two objects in any 2-category, preserved by any 2-functor:\\
1.4.6. Definition (equivalence). An equivalence in a 2-category is given by

\begin{itemize}
  \item[-] a pair of objects $A$ and $B$;
  \item[-] a pair of 1-cells $f: A \rightarrow B$ and $g: B \rightarrow A$; and
  \item[-] a pair of invertible 2-cells\\
\includegraphics[max width=\textwidth, alt={}, center]{b375083c-8174-49a1-a017-e977744038d8-33_124_256_1210_605}\\
and\\
\includegraphics[max width=\textwidth, alt={}, center]{b375083c-8174-49a1-a017-e977744038d8-33_121_242_1179_1274}
\end{itemize}

When $A$ and $B$ are equivalent, we write $A \simeq B$ and refer to the 1-cells $f$ and $g$ as equivalences, denoted by "\~{}."

In the case of the homotopy 2-category of an $\infty$-cosmos we have a competing definition of equivalence from 1.2.1: namely a 1-cell $f: A \xrightarrow{\leadsto} B$ that induces an equivalence $f_{*}: \operatorname{Fun}(X, A) \xrightarrow{\leadsto} \operatorname{Fun}(X, B)$ on functor spaces - or equivalently, by Lemma 1.2.15, a homotopy equivalence defined relative to the interval II. Crucially, all three notions of equivalence coincide:\\
1.4.7. Theorem (equivalences are equivalences). In any $\infty$-cosmos $\mathcal{K}$, the following are equivalent and characterize what it means for a functor $f: A \rightarrow B$ between $\infty$-categories to define an equivalence.

\begin{itemize}
  \item[(i)] For all $X \in \mathcal{K}$, the post-composition map $f_{*}: \operatorname{Fun}(X, A) \xrightarrow{\sim} \operatorname{Fun}(X, B)$ defines an equivalence of quasi-categories.
  \item[(ii)] There exists a functor $g: B \rightarrow A$ and natural isomorphisms $\alpha: \mathrm{id}_{A} \cong g f$ and $\beta: f g \cong \mathrm{id}_{B}$ in the homotopy 2-category.
  \item[(iii)] There exists a functor $g: B \rightarrow A$ and maps\\
\includegraphics[max width=\textwidth, alt={}, center]{b375083c-8174-49a1-a017-e977744038d8-33_312_245_2066_682}\\
and in the $\infty-\operatorname{cosmos} \mathcal{K}$.\\
\includegraphics[max width=\textwidth, alt={}, center]{b375083c-8174-49a1-a017-e977744038d8-33_311_236_2067_1212}
\end{itemize}

As an illustrative comparison of 2-categorical and quasi-categorical techniques, rather than appealing to Lemma 1.2.15 to prove (i)⇔(iii), we re-prove it.

Proof. For (i)⇒(ii), if the induced map $f_{*}: \operatorname{Fun}(X, A) \xrightarrow{\sim} \operatorname{Fun}(X, B)$ defines an equivalence of quasi-categories then the functor $f_{*}: \mathrm{hFun}(X, A) \leadsto \mathrm{hFun}(X, B)$ defines an equivalence of categories, by Remark 1.1.24. In particular, the equivalence $f_{*}: \operatorname{hFun}(B, A) \xrightarrow{\sim} \mathrm{hFun}(B, B)$ is essentially surjective so there exists $g \in \operatorname{hFun}(B, A)$ and an isomorphism $\beta: f g \cong \operatorname{id}_{B} \in \operatorname{hFun}(B, B)$. Now since $f_{*}: \operatorname{hFun}(A, A) \xrightarrow{\beth} \mathrm{hFun}(A, B)$ is fully faithful, the isomorphism $\beta f: f g f \cong f \in \mathrm{hFun}(A, B)$ can be lifted to define an isomorphism $\alpha^{-1}: g f \cong \operatorname{id}_{A} \in \operatorname{hFun}(A, A)$. This defines the data of a 2-categorical equivalence in Definition 1.4.6.

To see that (ii)⇒(iii) recall from Lemma 1.4.3 that the natural isomorphisms $\alpha: \mathrm{id}_{A} \cong g f$ and $\beta: f g \cong \operatorname{id}_{B}$ in $\mathfrak{h} \mathcal{K}$ are represented by maps $\alpha: A \rightarrow A^{\mathbb{I}}$ and $\beta: B \rightarrow B^{\mathbb{I}}$ in $\mathcal{K}$ as in (1.2.16).

Finally, (iii)⇒(i) since $\operatorname{Fun}(X,-)$ carries the data of (iii) to the data of an equivalence of quasi-categories as in Definition 1.1.23. \(\square\)

It is hard to overstate the importance of Theorem 1.4.7 for the work that follows. The categorical constructions that we introduce for $\infty$-categories, $\infty$-functors, and $\infty$-natural transformations are invariant under 2-categorical equivalence in the homotopy 2-category and the universal properties we develop similarly characterize 2 -categorical equivalence classes of $\infty$-categories. Theorem 1.4.7 then asserts that such constructions are "homotopically correct": both invariant under equivalence in the $\infty$-cosmos and precisely identifying equivalence classes of objects.

The equivalence invariance of the functor space in the codomain variable is axiomatic, but equivalence invariance in the domain variable is not. ${ }^{23}$ Nor is it evident how this could be proven from either (i) or (iii) of Theorem 1.4.7. But using (ii) and 2-categorical techniques, there is now a short proof.\\
1.4.8. Corollary. Equivalences of $\infty$-categories $A^{\prime} \xrightarrow{\sim} A$ and $B \xrightarrow{ } B^{\prime}$ induce an equivalence of functor spaces $\operatorname{Fun}(A, B) \xrightarrow{\sim} \operatorname{Fun}\left(A^{\prime}, B^{\prime}\right)$.

Proof. The representable simplicial functors $\operatorname{Fun}(A,-): \mathcal{K} \rightarrow Q$ Cat and $\operatorname{Fun}(-, B): \mathcal{K}^{\text {op }} \rightarrow$ QCat induce 2-functors $\operatorname{Fun}(A,-): \mathfrak{h K} \rightarrow \mathfrak{h Q C a t}$ and $\operatorname{Fun}(-, B): \mathfrak{h} \mathcal{K}^{\mathrm{op}} \rightarrow \mathfrak{h Q C a t}$, which preserve the 2-categorical equivalences of Definition 1.4.6. By Theorem 1.4.7 this is what we wanted to show. \(\square\)

There is also a standard 2-categorical notion of an isofibration, defined in the statement of Proposition 1.4.9 and elaborated upon in Definition B.4.4. We now show that any isofibration in an $\infty$-cosmos defines an isofibration in its homotopy 2-category.\\
1.4.9. Proposition (isofibrations are isofibrations). An isofibration $p: E \rightarrow B$ in an $\infty$-cosmos $\mathcal{K}$ also defines an isofibration in the homotopy 2-category $\mathfrak{h K}$ : given any invertible 2-cell as displayed below-left abutting to $B$ with a specified lift of one of its boundary 1-cells through $p$, there exists an invertible 2-cell abutting to $E$ with this boundary 1-cell as displayed below-right that whiskers with p to the original 2-cell.\\
\includegraphics[max width=\textwidth, alt={}, center]{b375083c-8174-49a1-a017-e977744038d8-34_209_770_2068_681}

\footnotetext{${ }^{23}$ Lemma 1.3.2 does not apply since Fun(-, $B$ ) is not cosmological.
}Proof. The universal property of the statement says that the functor

$$
p_{*}: \operatorname{hFun}(X, E) \rightarrow \operatorname{hFun}(X, B)
$$

is an isofibration of categories in the sense defined in Proposition 1.2.11. By axiom 1.2.1(ii), since $p: E \rightarrow B$ is an isofibration in $\mathcal{K}$, the induced map $p_{*}: \operatorname{Fun}(X, E) \rightarrow \operatorname{Fun}(X, B)$ is an isofibration of quasi-categories. So it suffices to show that the functor $\mathrm{h}: Q$ Cat $\rightarrow$ Cat carries isofibrations of quasi-categories to isofibrations of categories.

So let us now consider an isofibration $p: E \rightarrow B$ between quasi-categories. By Corollary 1.1.16, every isomorphism $\beta$ in the homotopy category $\mathrm{h} B$ of the quasi-category $B$ is represented by a simplicial $\operatorname{map} \beta: \mathbb{I} \rightarrow B$. By Definition 1.1.17, the lifting problem\\
\includegraphics[max width=\textwidth, alt={}, center]{b375083c-8174-49a1-a017-e977744038d8-35_204_203_741_964}\\
can be solved, and the map $\gamma: \mathbb{I} \rightarrow E$ so produced represents a lift of the isomorphism from $\mathrm{h} B$ to an isomorphism in $\mathrm{h} E$ with domain $e$. \(\square\)\\
1.4.10. Convention (on isofibrations in homotopy 2-categories). Since the converse to Proposition 1.4.9 does not hold, there is a potential ambiguity when using the term "isofibration" to refer to a map in the homotopy 2-category of an $\infty$-cosmos. We adopt the convention that when we declare a map in $\mathfrak{h} \mathcal{K}$ to be an isofibration we always mean this is the stronger sense of defining an isofibration in $\mathcal{K}$. This stronger condition gives us access to the 2-categorical lifting property of Proposition 1.4.9 and also to homotopical properties axiomatized in Definition 1.2.1, which ensure that the strictly defined limits of 1.2.1(i) are automatically equivalence invariant constructions (see §C. 1 and Proposition 6.2.8).

We conclude this chapter with a final definition that can be extracted from the homotopy 2-category of an $\infty$-cosmos. The 1- and 2-cells in the homotopy 2-category from the terminal $\infty$-category $1 \in \mathcal{K}$ to a generic $\infty$-category $A \in \mathcal{K}$ define the objects and morphisms in the homotopy category of the $\infty$-category $A$.\\
1.4.11. Definition (homotopy category of an $\infty$-category). The homotopy category of an $\infty$-category $A$ in an $\infty$-cosmos $\mathcal{K}$ is defined to be the homotopy category of its underlying quasi-category, that is:

$$
\mathrm{h} A:=\mathrm{hFun}(1, A):=\mathrm{h}(\operatorname{Fun}(1, A)) .
$$

As we shall discover, homotopy categories generally inherit "derived" analogues of structures present at the level of $\infty$-categories. An early example of this appears in Proposition 2.1.7(ii).

\section*{Exercises.}
1.4.i. Exercise.

\begin{itemize}
  \item[(i)] What is the homotopy 2-category of the $\infty$-cosmos Cat of 1-categories?
  \item[(ii)] Prove that the nerve defines a 2-functor $C a t \hookrightarrow \mathfrak{h} Q C a t$ that is locally fully faithful.
\end{itemize}

1.4.ii. Exercise. Demonstrate that the homotopy 2-category of the dual $\operatorname{cosmos} \mathcal{K}^{\mathrm{co}}$ of an $\infty$-cosmos $\mathcal{K}$ is the co-dual of the homotopy 2-category $\mathfrak{h} \mathcal{K}$ - in symbols $\mathfrak{h}\left(\mathcal{K}^{\mathrm{C}^{\mathrm{O}}}\right) \cong(\mathfrak{h} \mathcal{K})^{\text {co }}$ - with the domains and codomains of 2-cells but not 1-cells reversed (see Definition B.1.6).\\
1.4.iii. Exercise. Consider a natural isomorphism $A \xrightarrow[g]{f} B$ between a parallel pair of functors in an $\infty$-cosmos. Give two proofs that if either $f$ or $g$ is an equivalence then both functors are, either by arguing entirely in the homotopy 2-category or by appealing to Lemma 1.4.3.\\
1.4.iv. Exercise. Extend Lemma 1.2.27 to show that the following four conditions are equivalent, characterizing the discrete objects $E$ in an $\infty$-cosmos $\mathcal{K}$ :

\begin{itemize}
  \item[(i)] $E$ is a discrete object in the homotopy 2-category $\mathfrak{h} \mathcal{K}$, that is, every 2-cell with codomain $E$ is invertible.
  \item[(ii)] For each $X \in \mathcal{K}$, the hom-category $\mathrm{hFun}(X, E)$ is a groupoid.
  \item[(iii)] For each $X \in \mathcal{K}$, the mapping quasi-category $\operatorname{Fun}(X, E)$ is a Kan complex.
  \item[(iv)] The isofibration $E^{\mathbb{I}} \rightarrow E^{2}$, induced by the inclusion of simplicial sets $2 \hookrightarrow \mathbb{I}$, is a trivial fibration.
\end{itemize}

1.4.v. Exercise (10.3.1). Extend Lemma 1.4.4 to show that if $F: \mathcal{K} \rightarrow \mathcal{L}$ is a cosmological biequivalence then $F: \mathfrak{h} \mathcal{K} \rightarrow \mathfrak{h} \mathcal{L}$ is a 2-categorical biequivalence, a 2-functor that is essentially surjective on objects up to equivalence that locally defines an equivalence of hom-categories.\\
\href{http://1.4.vi}{1.4.vi}. Exercise. Let $F: \mathcal{K} \xrightarrow{\sim} \mathcal{L}$ be a cosmological biequivalence and let $A, B \in \mathcal{K}$. Re-prove part of the statement of Lemma 1.3.12: that if $F A \simeq F B$ in $\mathcal{L}$ then $A \simeq B$ in $\mathcal{K}$.\\
1.4.vii. Exercise (3.6.2). Let $B$ be an $\infty$-category in the $\infty$-cosmos $\mathcal{K}$ and let $\mathfrak{h} \mathcal{K}_{/ B}$ denote the 2-category whose

\begin{itemize}
  \item[-] objects are isofibrations $E \rightarrow B$ in $\mathcal{K}$ with codomain $B$;
  \item[-] 1-cells are 1-cells in $\mathfrak{h} \mathcal{K}$ over $B$; and
\end{itemize}

\begin{center}
\includegraphics[max width=\textwidth, alt={}]{b375083c-8174-49a1-a017-e977744038d8-36_136_249_1455_938}
\end{center}

\begin{itemize}
  \item[-] 2-cells are 2-cells $\alpha$ in $\mathfrak{h} \mathcal{K}$\\
\includegraphics[max width=\textwidth, alt={}, center]{b375083c-8174-49a1-a017-e977744038d8-36_221_250_1661_937}\\
that lie over $B$ in the sense that $q \alpha=\mathrm{id}_{p}$.
\end{itemize}

Argue that the homotopy 2-category $\mathfrak{h}\left(\mathcal{K}_{/ B}\right)$ of the sliced $\infty$-cosmos has the same underlying 1-category but different 2-cells. How do these compare with the 2-cells of $\mathfrak{h} \mathcal{K}_{/ B}$ ?

\begin{itemize}
  \item[-] 
  \item[-] 
\end{itemize}


\end{document}